\documentclass[11pt,a4paper]{article}

\usepackage[T1]{fontenc}
\usepackage[utf8]{inputenc}
\usepackage[french]{babel}
\usepackage{lmodern}
\usepackage{microtype}
\usepackage[margin=2.4cm]{geometry}
\usepackage{amsmath,amssymb,amsthm,mathtools}
\usepackage{booktabs}
\usepackage{tabularx}
\usepackage{enumitem}
\usepackage{xcolor}
\usepackage{tikz}
\usetikzlibrary{arrows.meta,positioning,shapes.geometric,calc,fit,backgrounds}
\usepackage{algorithm}
\usepackage{algpseudocode}
\usepackage{listings}
\usepackage[hidelinks,colorlinks=true,linkcolor=Bleu,citecolor=Bleu,urlcolor=Bleu]{hyperref}

\definecolor{Bleu}{HTML}{1F4E79}
\definecolor{Gris}{HTML}{6B7280}
\definecolor{GrisClair}{HTML}{F1F2F4}
\definecolor{Rouge}{HTML}{A83232}
\definecolor{Vert}{HTML}{2E6B4F}

% --- Environnements de theoremes ---
\theoremstyle{definition}
\newtheorem{definition}{Définition}[section]
\newtheorem{propriete}{Propriété}[section]
\theoremstyle{remark}
\newtheorem{remarque}{Remarque}[section]
\newtheorem{avertissement}{Avertissement}[section]

% --- Algorithmes en francais ---
\floatname{algorithm}{Algorithme}
\renewcommand{\listalgorithmname}{Liste des algorithmes}
\algrenewcommand\algorithmicrequire{\textbf{Entrées :}}
\algrenewcommand\algorithmicensure{\textbf{Sortie :}}
\algrenewcommand\algorithmicfor{\textbf{pour}}
\algrenewcommand\algorithmicforall{\textbf{pour tout}}
\algrenewcommand\algorithmicdo{\textbf{faire}}
\algrenewcommand\algorithmicwhile{\textbf{tant que}}
\algrenewcommand\algorithmicend{\textbf{fin}}
\algrenewcommand\algorithmicif{\textbf{si}}
\algrenewcommand\algorithmicthen{\textbf{alors}}
\algrenewcommand\algorithmicelse{\textbf{sinon}}
\algrenewcommand\algorithmicreturn{\textbf{retourner}}

% --- Listings Python ---
\lstset{
  language=Python,
  basicstyle=\ttfamily\footnotesize,
  keywordstyle=\color{Bleu}\bfseries,
  commentstyle=\color{Vert}\itshape,
  stringstyle=\color{Rouge},
  backgroundcolor=\color{GrisClair},
  frame=single,
  rulecolor=\color{Gris},
  breaklines=true,
  showstringspaces=false,
  columns=fullflexible,
  extendedchars=true,
  literate=%
    {é}{{\'e}}1 {è}{{\`e}}1 {ê}{{\^e}}1 {à}{{\`a}}1 {ù}{{\`u}}1
    {ç}{{\c c}}1 {î}{{\^i}}1 {ô}{{\^o}}1 {û}{{\^u}}1 {É}{{\'E}}1
    {â}{{\^a}}1 {ë}{{\"e}}1 {ï}{{\"i}}1 {°}{{$^\circ$}}1 {œ}{{\oe}}1
}

% --- Raccourcis mathematiques ---
\newcommand{\R}{\mathbb{R}}
\newcommand{\N}{\mathbb{N}}
\newcommand{\Z}{\mathbb{Z}}
\newcommand{\PP}{\mathbb{P}}
\newcommand{\EE}{\mathbb{E}}
\newcommand{\Unif}{\mathcal{U}}
\newcommand{\Bball}{\mathbb{B}}
\newcommand{\Sphere}{\mathbb{S}}
\newcommand{\norm}[1]{\left\lVert #1 \right\rVert}
\newcommand{\abs}[1]{\left\lvert #1 \right\rvert}
\newcommand{\scal}[2]{\left\langle #1 ,\, #2 \right\rangle}
\newcommand{\ind}[1]{\mathbf{1}\{#1\}}
\newcommand{\simplex}{\Delta^{d-1}}
\newcommand{\Tmap}{\mathbf{T}}
\newcommand{\pareto}{\succ_{\mathrm{P}}}

\title{\textbf{Agrégation empirique de métriques de vulnérabilité}\\[0.3em]
\large Construction, estimation et validation d'un score synthétique\\ sur les produits de la nomenclature combinée}
\author{Note méthodologique}
\date{\today}

\begin{document}
\maketitle

\begin{abstract}
\noindent
Ce document recense les méthodes permettant d'ordonner un ensemble de vecteurs de $\R^d$ (par exemple dont les composantes représentent des métriques de vulnérabilité calculées pour chaque nomenclature), sous la contrainte que \emph{tous} les
coefficients de la synthèse soient déterminés empiriquement, sans hiérarchisation \emph{a priori}
des composantes des vecteurs. Trois familles de méthodes sont explorées : les ordres partiels sans aucun poids (dominance de
Pareto), les pondérations endogènes (entropie, CRITIC, ACP, \emph{benefit of the doubt}) associées
à différentes fonctions d'agrégation, et les approches géométriques multivariées, dont les rangs
de Kantorovitch fondés.
Une dernière section propose un protocole de comparaison et de contrôle de cohérence entre
agrégations.
\end{abstract}

\vspace{0.5em}
{\small\tableofcontents}
\newpage

%=======================================================================
\section{Cadrage et notations}
%=======================================================================

\subsection{Le problème}

\medskip

On dispose de $n$ produits de la nomenclature combinée, indexés par $i \in [n] := \{1,\dots,n\}$,
et de $d$ métriques de vulnérabilité (indice de Herfindahl--Hirschman de concentration des
fournisseurs, indices de dépendance, etc.), indexées par $j \in [d]$. Chaque produit est décrit
par un vecteur ligne $x_i = (x_{i1},\dots,x_{id}) \in \R^d$, et l'ensemble des observations forme
la matrice
\[
X = (x_{ij})_{i\in[n],\,j\in[d]} \in \R^{n\times d}.
\]
L'objectif est de construire une application de score $s : \R^d \to \R$ (ou, plus généralement,
un ordre sur $\{x_1,\dots,x_n\}$) permettant de désigner les produits les plus vulnérables, en
imposant que tout paramètre de $s$ soit \emph{estimé à partir de $X$} et non fixé par jugement
d'expert.

\bigskip

\subsection{L'obstacle théorique}

\medskip

Il n'existe pas d'ordre canonique sur $\R^d$ pour $d \ge 2$ : aucun ordre total ne peut
simultanément être compatible avec la structure d'espace vectoriel et invariant par permutation
des coordonnées. Tout classement linéaire des produits repose donc nécessairement sur une
hypothèse supplémentaire. Refuser de fixer des poids \emph{a priori} ne fait pas disparaître cette
hypothèse : cela la déplace vers une propriété statistique des données, qui joue alors le rôle du
jugement de valeur.

\medskip

\begin{avertissement}[Le coût caché de la pondération endogène]
  \label{av:cout}
  Chaque méthode « objective » substitue au jugement d'expert un postulat implicite :
  \begin{itemize}[nosep,leftmargin=1.4em]
    \item \textbf{entropie, CRITIC, ACP} : \emph{l'information utile est la variance}. Une métrique
    qui discrimine peu entre produits reçoit un poids faible --- même si elle capture le mécanisme
    de vulnérabilité le plus grave. Si toutes les positions NC sont exposées de la même manière à un
    risque majeur, ce risque sera mécaniquement neutralisé par le score.
    \item \textbf{ACP}: \emph{la vulnérabilité est la direction de plus grande dispersion}. Rien ne
    garantit que le premier axe factoriel soit un axe de vulnérabilité plutôt que, par exemple, un
    axe de volume d'échange ou un artefact d'échelle.
    \item \textbf{\emph{benefit of the doubt}}: \emph{chaque produit doit être jugé sous la
    pondération qui lui est la plus favorable}, ce qui, pour un indice de vulnérabilité, revient à
    une lecture de précaution (voir §\ref{sec:bod}).
    \item \textbf{transport optimal}: \emph{la position dans le nuage résume l'information}, avec
    une notion de centralité \emph{multivariée} qui absorbe corrélations et échelles.
  \end{itemize}
\end{avertissement}

\subsection{Vocabulaire}

\begin{definition}[Compensation]
  Une fonction d'agrégation $s$ est dite \emph{compensatoire} si une dégradation sur une coordonnée
  peut être intégralement compensée par une amélioration sur une autre. Par exemple :
  medskip
  \begin{itemize}
    \item La somme pondérée est totalement compensatoire ; 
    \item l'ordre lexicographique ne est pas du tout compensatoire ;
    \item la moyenne géométrique et les indices à pénalité occupent des positions intermédiaires.
  \end{itemize}
\end{definition}

\bigskip

\begin{definition}[Polarité]
  Une métrique est de polarité \emph{positive} si une valeur élevée signale une vulnérabilité
  élevée, et de polarité \emph{négative} dans le cas contraire.
\end{definition}

\bigskip

\begin{definition}[Dominance de Pareto]
  \label{def:pareto}
  Après orientation de toutes les métriques en polarité positive, le produit $i$ \emph{domine} le
  produit $k$, noté $x_i \pareto x_k$, si :

  \medskip

  \[
  \forall j \in [d],\ x_{ij} \ge x_{kj}
  \quad\text{et}\quad
  \exists j_0 \in [d],\ x_{ij_0} > x_{kj_0}.
  \]
\end{definition}

\bigskip

\begin{propriete}[Monotonie, ou compatibilité paretienne]
  \label{prop:monotonie}
  Une fonction de score $s$ est \emph{compatible avec la dominance} si
  $x_i \pareto x_k \Rightarrow s(x_i) > s(x_k)$.
\end{propriete}

La propriété \ref{prop:monotonie} est le seul axiome que l'on peut exiger sans engager de
jugement de valeur. Elle fournit à ce titre le test de cohérence le plus discriminant du
protocole de validation (§\ref{sec:violation}).

\subsection{Présentation synthétique des méthodes}

\medskip

\begin{table}[htbp]
  \centering
  \small
  \renewcommand{\arraystretch}{1.25}
  \begin{tabularx}{\textwidth}{@{}l l c c X@{}}
  \toprule
  \textbf{Méthode} & \textbf{§} & \textbf{Poids} & \textbf{Compens.} & \textbf{Postulat implicite} \\
  \midrule
  Dominance de Pareto        & \ref{sec:pareto} & aucun      & non      & aucun (ordre partiel) \\
  Tri non dominé             & \ref{sec:tri}    & aucun      & non      & les fronts successifs sont ordonnés \\
  Entropie de Shannon        & \ref{sec:entropie} & endogène & variable & information $=$ dispersion \\
  CRITIC                     & \ref{sec:critic} & endogène   & variable & information $=$ dispersion $\times$ originalité \\
  ACP / axe 1                & \ref{sec:acp}    & endogène   & oui      & vulnérabilité $=$ axe de plus grande variance \\
  \emph{Benefit of the doubt}& \ref{sec:bod}    & endogène, individualisée & oui & jugement sous la pondération la plus favorable \\
  Moyenne géométrique        & \ref{sec:geo}    & endogène   & partielle& les profils déséquilibrés sont pénalisés \\
  Indice MPI                 & \ref{sec:mpi}    & aucun      & partielle& idem, avec pénalité explicite \\
  TOPSIS                     & \ref{sec:topsis} & endogène   & oui      & proximité à un pôle idéal empirique \\
  Mahalanobis                & \ref{sec:maha}   & endogène   & oui      & la corrélation mesure la redondance \\
  Rang de Kantorovitch       & \ref{sec:ot}     & endogène   & oui      & vulnérabilité $=$ position dans le nuage \\
  SMAA                       & \ref{sec:smaa}   & aléatoires & --       & ignorance $=$ loi uniforme sur le simplexe \\
  \bottomrule
  \end{tabularx}
  \caption{Vue d'ensemble. La colonne « Poids » indique si la méthode requiert des coefficients et,
  le cas échéant, comment ils sont obtenus.}
\end{table}

%=======================================================================
\section{Préalables : polarité, normalisation, diagnostic}
\label{sec:prealables}
%=======================================================================

\medskip

Le choix de normalisation influence le
classement final ; il doit
être traité comme une décision méthodologique à part entière et versé à l'analyse de sensibilité.

\bigskip

\subsection{Orientation des métriques}

\medskip

Toutes les coordonnées doivent être mises en polarité positive. Pour une métrique de polarité
négative, on applique $x_{ij} \mapsto -x_{ij}$, ou $x_{ij} \mapsto \max_k x_{kj} - x_{ij}$, ou
encore $x_{ij}\mapsto 1/x_{ij}$ si la métrique est strictement positive et que l'inversion a un
sens économique. Ces trois transformations ne sont pas équivalentes : seule la première est
affine, donc seule elle laisse invariants les résultats des méthodes fondées sur la covariance
(ACP, Mahalanobis, transport optimal dans le cas elliptique).

\bigskip

\subsection{Traitement des queues de distribution}

\medskip

Les indices de concentration (HHI et apparentés) produisent typiquement des distributions très
asymétriques : une poignée de positions NC atteignent des valeurs extrêmes tandis que la masse est
concentrée près de $0$. Trois attitudes possibles, à documenter :

\medskip

\begin{enumerate}[nosep,leftmargin=1.6em]
  \item \textbf{Ne rien faire.} Les extrêmes sont précisément les produits que l'on cherche à
  identifier. Cohérent, mais la normalisation min--max sera pilotée par un unique produit, et
  l'ACP par quelques points de levier.
  \item \textbf{Winsorisation} au quantile $q$ (typiquement $q = 0{,}99$) : $x_{ij} \mapsto
  \min(x_{ij}, F_j^{-1}(q))$. Stabilise les statistiques d'échelle sans écraser l'ordre.
  \item \textbf{Transformation monotone} ($\log(1+x)$, racine, transformation de Box--Cox
  ajustée par maximum de vraisemblance). Le paramètre de Box--Cox est estimé sur les données.
\end{enumerate}

\medskip

\begin{remarque}
  Une transformation strictement croissante appliquée coordonnée par coordonnée laisse la relation
  de dominance $\pareto$ \emph{invariante}: le front de Pareto est donc insensible à ce choix.
  C'est un argument fort en faveur du calcul systématique de ce front comme point de référence
  (§\ref{sec:pareto}).
\end{remarque}

\bigskip

\subsection{Familles de normalisation}

\medskip

\begin{table}[htbp]
  \centering
  \small
  \renewcommand{\arraystretch}{1.4}
  \begin{tabularx}{\textwidth}{@{}l l X@{}}
    \toprule
    \textbf{Schéma} & \textbf{Formule} & \textbf{Effet sur le classement} \\
    \midrule
    Min--max &
    $\tilde{x}_{ij} = \dfrac{x_{ij}-\min_k x_{kj}}{\max_k x_{kj}-\min_k x_{kj}}$ &
    Borne dans $[0,1]$, conserve les écarts relatifs. Très sensible aux deux valeurs extrêmes:
    un seul produit atypique comprime toute la métrique. \\
    Centrage--réduction &
    $\tilde{x}_{ij} = \dfrac{x_{ij}-\bar{x}_j}{\sigma_j}$ &
    Égalise les variances, donc \emph{égalise implicitement les poids} dans une somme. Non borné ;
    conserve l'asymétrie. \\
    Robuste &
    $\tilde{x}_{ij} = \dfrac{x_{ij}-\mathrm{med}_j}{\mathrm{MAD}_j}$ &
    Idem, insensible aux extrêmes. Recommandé avec des HHI. \\
    Rangs &
    $\tilde{x}_{ij} = \dfrac{\mathrm{rg}(x_{ij})-1}{n-1}$ &
    Détruit toute information cardinale (un écart énorme et un écart infime deviennent identiques),
    mais rend le score invariant à toute transformation monotone. \\
    Quantile gaussien &
    $\tilde{x}_{ij} = \Phi^{-1}\!\big(\tfrac{\mathrm{rg}(x_{ij})-1/2}{n}\big)$ &
    Marges gaussiennes, corrélations préservées en structure (copule). Utile en amont d'une ACP ou
    d'un transport optimal. \\
    \bottomrule
  \end{tabularx}
  \caption{Schémas de normalisation. Aucun n'est neutre; le centrage--réduction en particulier
  constitue déjà une pondération déguisée.}
  \label{tab:normalisation}
\end{table}

\bigskip

\subsection{Diagnostic de corrélation}

\medskip

Avant toute agrégation, examiner la matrice de corrélation $R = (r_{jk})$ des métriques
normalisées, de préférence par corrélation de Spearman (monotone, robuste). Ce diagnostic
conditionne le choix de la famille de méthodes :

\medskip

\begin{itemize}[nosep,leftmargin=1.4em]
  \item \textbf{Métriques fortement corrélées} ($\abs{r_{jk}}$ élevés): une somme pondérée compte
  plusieurs fois la même information sous-jacente. Privilégier ACP, Mahalanobis, CRITIC ou
  transport optimal, qui décorrèlent explicitement.
  \item \textbf{Métriques quasi indépendantes}: l'ACP n'a pas d'axe dominant, son premier axe est
  arbitraire et instable. Privilégier les pondérations par dispersion ou les ordres partiels.
  \item \textbf{Corrélations négatives}: signale des mécanismes de vulnérabilité antagonistes.
  L'agrégation en un score unique perd alors beaucoup d'information; envisager plusieurs scores
  thématiques, ou une méthode acceptant l'incomparabilité.
\end{itemize}

\bigskip

On peut compléter cette première analyse par le calcul de l'indice KMO et le test de sphéricité de Bartlett pour évaluer si une structure
factorielle est seulement plausible.

\bigskip

\begin{algorithm}[htbp]
  \caption{Prétraitement}
  \label{algo:pretraitement}
  \begin{algorithmic}[1]
    \Require matrice brute $X\in\R^{n\times d}$, vecteur de polarités $p\in\{-1,+1\}^d$,
    quantile de winsorisation $q$, schéma de normalisation $\nu$
    \Ensure matrice normalisée $\tilde{X}$, matrice de corrélation $R$
    \For{$j = 1$ \textbf{à} $d$}
      \If{$p_j = -1$} $x_{\cdot j} \gets -x_{\cdot j}$ \Comment{Orientation en polarité positive}
      \EndIf
      \State $x_{\cdot j} \gets \min\!\big(x_{\cdot j},\, \hat{F}_j^{-1}(q)\big)$
            \Comment{Winsorisation des valeurs extrêmes}
      \State $\tilde{x}_{\cdot j} \gets \nu(x_{\cdot j})$
            \Comment{Application du schéma retenu, cf. table~\ref{tab:normalisation}}
    \EndFor
    \State $R \gets \operatorname{corr}_{\text{Spearman}}(\tilde{X})$
    \State \textbf{Contrôle :} valeurs manquantes, colonnes constantes, produits dupliqués
    \State \Return $\tilde{X}$, $R$
  \end{algorithmic}
\end{algorithm}

\bigskip

\begin{avertissement}
  Les positions NC auxuqelles sont associées de très faibles volumes produisent des métriques de concentration numériquement
  extrêmes mais statistiquement non significatives (un HHI calculé sur trois déclarations vaut
  mécaniquement près de $1$). Il est préférable de \emph{filtrer explicitement} ces positions,
  ou de leur associer un intervalle de confiance et de propager cette incertitude dans le
  classement (§\ref{sec:bootstrap}), plutôt que de laisser le score les hisser en tête.
\end{avertissement}

\bigskip

%=======================================================================
\section{Ordres partiels : classer sans aucun coefficient}
\label{sec:pareto}
%=======================================================================

\subsection{Idée}

C'est la seule construction totalement dépourvue d'hypothèse. On n'ordonne que les paires de
produits pour lesquelles \emph{toutes} les métriques concordent (définition \ref{def:pareto}).
Le résultat n'est pas un classement complet mais un ordre partiel, dont l'élément saillant est
le \emph{front de Pareto} :

\medskip

\[
\mathcal{F}_1 = \big\{ i \in [n] \;:\; \nexists\, k \in [n],\ x_k \pareto x_i \big\},
\]

\bigskip

l'ensemble des produits non dominés, c'est-à-dire ceux que l'on ne peut écarter du haut du
classement sans déclarer arbitrairement qu'un critère prime sur un autre.

\bigskip

\begin{figure}[htbp]
  \centering
  \begin{tikzpicture}[scale=1.0]
    \draw[->,thick,Gris] (0,0) -- (7.2,0) node[right,black,font=\small] {métrique 1};
    \draw[->,thick,Gris] (0,0) -- (0,5.0) node[above,black,font=\small] {métrique 2};
    % points domines
    \foreach \x/\y in {0.8/1.0, 1.5/0.6, 2.0/1.9, 1.2/2.4, 2.9/1.2, 3.5/2.2, 2.4/3.0,
                      4.2/1.0, 4.9/1.8, 3.1/0.5, 5.4/1.1, 2.6/2.2, 0.6/3.0, 4.0/2.9}
      \fill[Gris] (\x,\y) circle (2.1pt);
    % front de Pareto
    \foreach \x/\y in {0.9/4.4, 2.2/4.1, 3.6/3.7, 5.0/3.3, 6.1/2.4, 6.6/1.4}
      \fill[Rouge] (\x,\y) circle (3pt);
    \draw[Rouge,thick,dashed] (0.9,4.4) -- (2.2,4.1) -- (3.6,3.7) -- (5.0,3.3) -- (6.1,2.4) -- (6.6,1.4);
    \node[Rouge,font=\small\bfseries] at (3.0,4.65) {front de Pareto $\mathcal{F}_1$};
    % cone de dominance
    \begin{scope}[opacity=0.18]
      \fill[Bleu] (3.5,2.2) rectangle (7.0,4.8);
    \end{scope}
    \draw[Bleu,thick,-{Stealth}] (3.5,2.2) -- (4.6,2.2);
    \draw[Bleu,thick,-{Stealth}] (3.5,2.2) -- (3.5,3.3);
    \node[Bleu,font=\scriptsize,align=center] at (5.7,4.15) {cône des points\\ qui dominent $x_i$};
    \fill[Bleu] (3.5,2.2) circle (2.5pt) node[below left,font=\scriptsize,black] {$x_i$};
  \end{tikzpicture}
  \caption{Dominance de Pareto en dimension $2$, les deux métriques étant orientées en polarité
  positive. Un produit est dominé s'il existe au moins un point dans le cône supérieur droit issu
  de lui.}
\end{figure}

\bigskip

\subsection{Tri non dominé et raffinements ordinaux}
\label{sec:tri}

\medskip

Le front seul ne discrimine pas à l'intérieur de l'ensemble qu'il détermine, aussi trois raffinements, tous sans pondération, peuvent leur être associés:

\medskip

\paragraph{Tri non dominé par couches.} On retire $\mathcal{F}_1$, on recalcule le front sur le
reste pour obtenir $\mathcal{F}_2$, et ainsi de suite. Le rang de couche $\ell(i)$ définit un
classement ordinal grossier. C'est la procédure de tri de NSGA-II (algorithme \ref{algo:nds}), de
complexité $O(d\,n^2)$ dans sa version naïve.

\bigskip

\paragraph{Comptage de dominance.} Pour chaque produit,
\[
\delta(i) = \underbrace{\#\{k : x_i \pareto x_k\}}_{\text{ce qu'il domine}}
          - \underbrace{\#\{k : x_k \pareto x_i\}}_{\text{ce qui le domine}},
\]
qui fournit un score entier plus finement discriminant que la couche. C'est un raffinement de
l'ordre partiel : si $x_i \pareto x_k$ alors $\delta(i) > \delta(k)$, la
propriété \ref{prop:monotonie} est donc satisfaite par construction.

\bigskip

\paragraph{Profondeur de dominance normalisée.} $\delta(i)/(n-1) \in [-1,1]$, directement
comparable entre jeux de données de tailles différentes.

\medskip

\begin{algorithm}[htbp]
  \caption{Tri non dominé et comptage de dominance}
  \label{algo:nds}
  \begin{algorithmic}[1]
  \Require $\tilde{X}\in\R^{n\times d}$ (métriques orientées en polarité positive)
  \Ensure couches $(\mathcal{F}_\ell)_\ell$, comptage $\delta \in \Z^n$
  \For{$i = 1$ \textbf{à} $n$}
    \State $S_i \gets \emptyset$ \Comment{Ensemble des produits dominés par $i$}
    \State $c_i \gets 0$ \Comment{Nombre de produits dominant $i$}
    \For{$k = 1$ \textbf{à} $n$, $k \neq i$}
      \If{$\tilde{x}_i \pareto \tilde{x}_k$} $S_i \gets S_i \cup \{k\}$
      \ElsIf{$\tilde{x}_k \pareto \tilde{x}_i$} $c_i \gets c_i + 1$
      \EndIf
    \EndFor
    \State $\delta_i \gets \abs{S_i} - c_i$
  \EndFor
  \State $\ell \gets 1$~;\quad $\mathcal{F}_1 \gets \{i : c_i = 0\}$
  \While{$\mathcal{F}_\ell \neq \emptyset$}
    \State $\mathcal{Q} \gets \emptyset$
    \ForAll{$i \in \mathcal{F}_\ell$}
      \ForAll{$k \in S_i$}
        \State $c_k \gets c_k - 1$
        \If{$c_k = 0$} $\mathcal{Q} \gets \mathcal{Q}\cup\{k\}$ \Comment{Le produit affleure la couche suivante}
        \EndIf
      \EndFor
    \EndFor
    \State $\ell \gets \ell + 1$~;\quad $\mathcal{F}_\ell \gets \mathcal{Q}$
  \EndWhile
  \State \Return $(\mathcal{F}_\ell)_\ell$, $\delta$
  \end{algorithmic}
\end{algorithm}

\bigskip

\subsection{Limites et stratégies de résolution}

\medskip

La proportion attendue de points non dominés croît très vite avec $d$ : pour $n$ points de
coordonnées indépendantes, l'espérance du cardinal du front est de l'ordre de
$(\ln n)^{d-1}/(d-1)!$. Avec $d = 3$ et $n = 5\,000$, cela reste modeste ; avec $d \ge 8$ le front
absorbe une part importante de la nomenclature et cesse d'être informatif. Il existe deux moyens de pallier ce comportement :

\medskip

\begin{itemize}[nosep,leftmargin=1.4em]
  \item \textbf{Cône élargi} ($\varepsilon$-dominance): on remplace $\pareto$ par
  $x_i \pareto^{\varepsilon} x_k \iff \forall j,\ x_{ij} \ge x_{kj} - \varepsilon_j$, ce qui
  élargit le cône de dominance et réduit le front. Le vecteur $\varepsilon$ peut être calibré
  empiriquement, par exemple comme l'écart-type du bruit d'estimation de chaque métrique (obtenu
  par bootstrap), ce qui lui donne une interprétation statistique plutôt qu'arbitraire.
  \item \textbf{Réduction préalable de $d$} par regroupement des métriques redondantes (classification
  ascendante hiérarchique sur $1-\abs{r_{jk}}$), le front étant ensuite calculé sur les
  représentants de groupes.
\end{itemize}

\bigskip

\begin{remarque}[Rôle dans le dispositif]
  Le front de Pareto ne suffit jamais à conclure, mais il constitue un \emph{invariant} précieux :
  il ne dépend ni de la normalisation monotone, ni des poids, ni de la fonction d'agrégation.
  Tout score synthétique qui ne placerait pas les éléments de $\mathcal{F}_1$ en haut du classement
  doit être considéré comme suspect (§\ref{sec:violation}).
\end{remarque}

\bigskip

%=======================================================================
\section{Pondérations endogènes}
\label{sec:poids}
%=======================================================================

\medskip

On cherche ici un vecteur de poids $w = (w_1,\dots,w_d)$ appartenant au simplexe :

\medskip

\[
\simplex = \Big\{ w \in \R^d_+ \;:\; \textstyle\sum_{j=1}^d w_j = 1 \Big\},
\]

\bigskip

entièrement déterminé par $X$. Les quatre méthodes ci-dessous diffèrent par la statistique de $X$
qu'elles convertissent en poids.

%-----------------------------------------------------------------------
\subsection{Pondération par l'entropie de Shannon}
\label{sec:entropie}

\paragraph{Idée.} Une métrique dont les valeurs sont quasi identiques sur tous les produits
n'aide pas à les distinguer : elle porte peu d'\emph{information} au sens de Shannon, et reçoit
donc un poids faible. Une métrique dont les valeurs sont très étalées reçoit un poids fort.

\paragraph{Construction.} On normalise chaque colonne (positive) en distribution de probabilité,
\[
p_{ij} = \frac{\tilde{x}_{ij}}{\sum_{k=1}^n \tilde{x}_{kj}},
\]
puis on calcule l'entropie normalisée de la colonne $j$ et son degré de divergence:
\[
E_j = -\frac{1}{\ln n}\sum_{i=1}^n p_{ij}\ln p_{ij} \in [0,1],
\qquad
g_j = 1 - E_j,
\qquad
w_j = \frac{g_j}{\sum_{k=1}^d g_k}.
\]
Par convention $p_{ij}\ln p_{ij} = 0$ lorsque $p_{ij}=0$. Le maximum $E_j = 1$ est atteint quand la
métrique est constante ; on a alors $w_j = 0$.

\bigskip

\begin{avertissement}
  La méthode exige $\tilde{x}_{ij} \ge 0$ et est \emph{non invariante par translation} : ajouter une
  constante à une colonne modifie ses poids. Elle doit donc impérativement être appliquée après une
  normalisation min--max, et le résultat rapporté conjointement avec le schéma de normalisation
  utilisé. C'est un des points où l'analyse de sensibilité de la section \ref{sec:comparaison} est
  la plus instructive.
\end{avertissement}

\bigskip

%-----------------------------------------------------------------------
\subsection{CRITIC}
\label{sec:critic}

\medskip

\paragraph{Idée.} L'entropie ne regarde qu'une colonne à la fois et ignore la redondance : deux
métriques quasi identiques reçoivent chacune un poids plein, et le mécanisme qu'elles capturent
pèse donc double. CRITIC (\emph{CRiteria Importance Through Intercriteria Correlation}) corrige
cela en pondérant chaque métrique par le produit de sa \emph{dispersion} et de son
\emph{originalité} vis-à-vis des autres.

\medskip

\paragraph{Construction.} Sur données min--max normalisées, avec $\sigma_j$ l'écart-type de la
colonne $j$ et $r_{jk}$ la corrélation entre colonnes $j$ et $k$ :
\[
C_j = \sigma_j \sum_{k=1}^{d} \big(1 - r_{jk}\big),
\qquad
w_j = \frac{C_j}{\sum_{m=1}^d C_m}.
\]
Le terme $\sum_k (1 - r_{jk})$ est maximal lorsque $j$ est anticorrélée à toutes les autres
métriques, minimal lorsqu'elle en est une réplique. Une variante robuste remplace $\sigma_j$ par
le MAD et $r$ par la corrélation de Spearman ; elle est préférable ici.

\medskip

\begin{algorithm}[htbp]
  \caption{Poids par entropie et par CRITIC}
  \label{algo:poids}
  \begin{algorithmic}[1]
    \Require $\tilde{X}\in[0,1]^{n\times d}$ (normalisation min--max, polarité positive)
    \Ensure $w^{\mathrm{ent}}, w^{\mathrm{cri}} \in \simplex$
    \Statex \textbf{-- Entropie --}
    \For{$j = 1$ \textbf{à} $d$}
      \State $p_{\cdot j} \gets \tilde{x}_{\cdot j} / \sum_i \tilde{x}_{ij}$
            \Comment{Conversion de la colonne en distribution}
      \State $E_j \gets -\frac{1}{\ln n}\sum_i p_{ij}\ln p_{ij}$ \Comment{Convention $0\ln 0 = 0$}
      \State $g_j \gets 1 - E_j$
    \EndFor
    \State $w^{\mathrm{ent}} \gets g / \sum_j g_j$
    \Statex \textbf{-- CRITIC --}
    \State $R \gets \operatorname{corr}(\tilde{X})$~;\quad $\sigma \gets \operatorname{sd}(\tilde{X})$
    \For{$j = 1$ \textbf{à} $d$}
      \State $C_j \gets \sigma_j \sum_{k} (1 - r_{jk})$
            \Comment{Dispersion pondérée par l'originalité du critère}
    \EndFor
    \State $w^{\mathrm{cri}} \gets C / \sum_j C_j$
    \State \Return $w^{\mathrm{ent}}, w^{\mathrm{cri}}$
  \end{algorithmic}
\end{algorithm}

\bigskip

%-----------------------------------------------------------------------
\subsection{Analyse en composantes principales}
\label{sec:acp}

\paragraph{Idée.} Si les métriques sont corrélées, c'est qu'elles mesurent en partie une même
grandeur latente. L'ACP identifie cette grandeur comme la direction de variance maximale et
l'utilise comme indice composite.

\paragraph{Construction.} Sur données centrées-réduites $Z$, on diagonalise la matrice de
corrélation $R = \frac{1}{n-1}Z^\top Z$ :
\[
R = \sum_{m=1}^{d}\lambda_m\, a_m a_m^\top,
\qquad \lambda_1 \ge \dots \ge \lambda_d \ge 0,
\qquad \norm{a_m} = 1.
\]
Le score du premier axe est $s^{\mathrm{ACP}}(x_i) = \scal{a_1}{z_i}$. Trois précautions
indispensables :
\begin{enumerate}[nosep,leftmargin=1.6em]
  \item \textbf{Signe.} $a_1$ n'est défini qu'au signe près. On fixe la convention en imposant
  $\sum_j a_{1j} > 0$, de sorte qu'un score élevé corresponde bien à des métriques élevées.
  \item \textbf{Positivité des saturations.} Si certaines composantes de $a_1$ sont négatives, le
  score \emph{viole la monotonie} (propriété \ref{prop:monotonie}) : augmenter une vulnérabilité
  peut faire baisser l'indice. C'est rédhibitoire pour un usage administratif. On rapporte alors
  la part de variance du premier axe et l'on envisage soit une contrainte de positivité (ACP non
  négative, \emph{NMF}), soit une autre méthode.
  \item \textbf{Part de variance.} Si $\lambda_1/d$ est faible (disons $< 0{,}5$), il n'existe pas
  de facteur dominant ; le premier axe n'est alors qu'une direction parmi d'autres, et son
  instabilité par rééchantillonnage sera visible au bootstrap.
\end{enumerate}

\bigskip

\paragraph{Variante multi-axes (procédure OCDE--JRC).} Plutôt que le seul premier axe, on retient
les $M$ composantes de valeur propre $\lambda_m > 1$, on leur applique une rotation varimax pour
obtenir des saturations $\ell_{jm}$ interprétables, puis on définit :

\medskip

\[
w_j \;\propto\; \max_{m \le M} \ell_{jm}^2 \cdot \frac{\lambda_m}{\sum_{m'\le M}\lambda_{m'}} .
\]

\bigskip

Chaque métrique est ainsi rattachée au facteur qu'elle sature le plus, et chaque facteur pèse à
proportion de la variance qu'il explique. C'est la procédure de référence du \emph{Handbook on
Constructing Composite Indicators}.

\bigskip

%-----------------------------------------------------------------------
\subsection{\emph{Benefit of the doubt} : une pondération par produit}
\label{sec:bod}

\paragraph{Idée.} Les méthodes précédentes cherchent \emph{un} vecteur de poids valable pour tous.
L'approche \emph{benefit of the doubt} (BoD) attribue à chaque produit la pondération qui lui est la
plus favorable, sous la seule contrainte que cette pondération, appliquée à tous les autres
produits, ne produise aucun score supérieur à $1$.

\paragraph{Construction.} Pour chaque produit $o \in [n]$, on résout le programme linéaire

\medskip

\[
\begin{aligned}
s^{\mathrm{BoD}}(x_o) \;=\; \max_{w^{(o)} \in \R^d_+} \quad & \sum_{j=1}^d w^{(o)}_j \tilde{x}_{oj} \\
\text{s.c.} \quad & \sum_{j=1}^d w^{(o)}_j \tilde{x}_{ij} \le 1, \quad \forall i \in [n], \\
& w^{(o)}_j \ge 0, \quad \forall j \in [d].
\end{aligned}
\tag{BoD}
\]

\bigskip

Le score obtenu appartient à $(0,1]$ et vaut $1$ si et seulement s'il existe une pondération
positive sous laquelle le produit $o$ est le meilleur de tous.

\bigskip

\paragraph{Interprétation pour un indice de vulnérabilité.} Le sens du « bénéfice du doute »
mérite d'être explicité, car il s'inverse par rapport à l'usage habituel. Dans un indice de
performance, BoD est bienveillant: il donne à chaque unité sa meilleure chance. Ici, l'objet
mesuré est une vulnérabilité: l'optimum retenu est donc celui qui rend le produit \emph{aussi
vulnérable que possible}. La lecture est celle de la précaution: \emph{un produit est signalé
dès lors qu'il existe au moins une pondération admissible sous laquelle il est en tête}. Cela permet de mettre en évidence des vulnérabilités en étant averse au risque.

\bigskip

\begin{propriete}[Lien avec la dominance]
  $s^{\mathrm{BoD}}(x_o) = 1$ implique que $x_o$ n'est dominé par aucun autre produit. La réciproque
  est fausse: un produit du front de Pareto situé dans une « concavité » du front peut obtenir un
  score strictement inférieur à $1$. Le score BoD est donc un raffinement \emph{convexe} du front de
  Pareto, et il satisfait la monotonie de la propriété~\ref{prop:monotonie}.
\end{propriete}

\begin{avertissement}[Solutions dégénérées]
  Sans contrainte supplémentaire, le programme (BoD) attribue fréquemment tout le poids à une seule
  métrique et ignore les autres, ce qui vide l'indice
  composite de son sens. Le remède standard consiste à ajouter des \emph{restrictions de poids}
  proportionnelles, elles-mêmes calibrées empiriquement :

  \medskip

  \[
  \underline{\pi} \;\le\; \frac{w^{(o)}_j \tilde{x}_{oj}}{\sum_{k} w^{(o)}_k \tilde{x}_{ok}} \;\le\; \overline{\pi},
  \qquad \forall j \in [d],
  \]

  \bigskip

  où $\underline{\pi} = \kappa/d$ et $\overline{\pi} = 1/(\kappa d)$ avec, par exemple,
  $\kappa = 1/2$: aucune métrique ne peut contribuer à moins de la moitié ni à plus du double de sa
  part uniforme. Ces bornes restent des paramètres, mais elles portent sur la \emph{structure} de la
  pondération et non sur la hiérarchie des critères ; elles se prêtent bien à un test de robustesse
  en faisant varier $\kappa$.
\end{avertissement}

\bigskip

\begin{algorithm}[htbp]
\caption{Score \emph{benefit of the doubt} avec restrictions de poids}
\label{algo:bod}
\begin{algorithmic}[1]
\Require $\tilde{X}\in[0,1]^{n\times d}$, paramètre de restriction $\kappa \in (0,1]$
\Ensure scores $s^{\mathrm{BoD}} \in (0,1]^n$, pondérations individuelles $(w^{(o)})_{o\in[n]}$
\For{$o = 1$ \textbf{à} $n$}
  \State Construire l'objectif $\max_w \; \tilde{x}_o^\top w$
  \State Contraintes de non-domination : $\tilde{X} w \le \mathbf{1}_n$
  \State Contraintes de part : $\kappa/d \le (w_j \tilde{x}_{oj})/(\tilde{x}_o^\top w) \le 1/(\kappa d)$
         \Comment{Linéarisation immédiate en multipliant par $\tilde{x}_o^\top w > 0$}
  \State $(s^{\mathrm{BoD}}_o, w^{(o)}) \gets$ résolution du programme linéaire
\EndFor
\State \Return $s^{\mathrm{BoD}}$, $(w^{(o)})_o$
\end{algorithmic}
\end{algorithm}

\bigskip

\begin{remarque}[Coût]
  Le programme est résolu $n$ fois, chacun avec $n$ contraintes. Avec $n \sim 5\,000$ positions NC,
  c'est faisable mais non trivia ~: on réduit fortement le coût en n'imposant les contraintes de
  non-domination que sur l'enveloppe convexe supérieure du nuage (les seules contraintes actives),
  identifiable au préalable par calcul de l'enveloppe convexe ou du front de Pareto.
\end{remarque}

\bigskip

%-----------------------------------------------------------------------
\subsection{Récapitulatif : que produisent ces pondérations ?}

\medskip

Les quatre méthodes livrent des vecteurs de poids qui n'ont aucune raison de coïncider. Il ne
faut pas y voir un défaut mais un résultat: l'écart entre eux mesure exactement l'indétermination
que la contrainte « pas de poids \emph{a priori} » laisse subsister. Le protocole de la
section~\ref{sec:comparaison} est conçu pour exploiter cet écart plutôt que pour le masquer par le
choix d'un gagnant.

\bigskip

%=======================================================================
\section{Fonctions d'agrégation}
\label{sec:agregation}
%=======================================================================

\medskip

Une fois les poids obtenus, il reste à choisir comment les combiner. Ce choix est indépendant du
précédent et porte sur une question différente : \emph{jusqu'où une métrique favorable peut-elle
racheter une métrique défavorable ?}

\medskip

\begin{figure}[htbp]
\centering
\begin{tikzpicture}[scale=1.6]
  \foreach \dx/\titre/\type in {0/{Somme pondérée}/lin, 3.4/{Moyenne géométrique}/geo, 6.8/{Indice à pénalité}/pen}{
    \begin{scope}[xshift=\dx cm]
      \draw[->,Gris] (0,0) -- (2.3,0) node[right,black,font=\scriptsize] {$m_1$};
      \draw[->,Gris] (0,0) -- (0,2.3) node[above,black,font=\scriptsize] {$m_2$};
      \node[font=\scriptsize\bfseries,black] at (1.15,2.6) {\titre};
    \end{scope}
  }
  % iso-lignes lineaires
  \begin{scope}
    \foreach \c in {0.8,1.4,2.0} \draw[Bleu,thick] (\c,0.05) -- (0.05,\c);
  \end{scope}
  % iso-lignes geometriques (hyperboles)
  \begin{scope}[xshift=3.4cm]
    \foreach \c in {0.16,0.49,1.0}
      \draw[Bleu,thick,domain=0.08:2.2,samples=80,variable=\x] plot ({\x},{min(\c/\x,2.2)});
  \end{scope}
  % iso-lignes penalisees (encore plus courbees)
  \begin{scope}[xshift=6.8cm]
    \foreach \c in {0.35,0.75,1.25}
      \draw[Bleu,thick,domain=0.16:2.2,samples=80,variable=\x] plot ({\x},{min((\c/\x)^1.6,2.2)});
  \end{scope}
  \node[font=\scriptsize,Gris,align=center] at (1.15,-0.55) {compensation\\totale};
  \node[font=\scriptsize,Gris,align=center] at (4.55,-0.55) {compensation\\partielle};
  \node[font=\scriptsize,Gris,align=center] at (7.95,-0.55) {compensation\\faible};
\end{tikzpicture}
\vspace{0.6em}
\caption{Courbes d'iso-score pour trois familles d'agrégation, avec deux métriques. Plus les
courbes sont incurvées vers l'origine, plus un profil déséquilibré est pénalisé à somme égale.}
\end{figure}

\bigskip

%-----------------------------------------------------------------------
\subsection{Somme pondérée}

\medskip

\[
s^{\mathrm{lin}}(x_i) = \sum_{j=1}^d w_j\,\tilde{x}_{ij}.
\]
Totalement compensatoire, monotone, immédiate à expliquer et à auditer. Son défaut est aussi sa
propriété : un produit médiocre partout et un produit catastrophique sur une seule dimension
peuvent obtenir le même score. Si la vulnérabilité opérationnelle relève d'un « maillon faible »,
cette agrégation est mal spécifiée.

%-----------------------------------------------------------------------
\subsection{Moyenne géométrique pondérée}
\label{sec:geo}

\medskip

\[
s^{\mathrm{geo}}(x_i) = \prod_{j=1}^d \tilde{x}_{ij}^{\,w_j}
= \exp\Big(\sum_{j=1}^d w_j \ln \tilde{x}_{ij}\Big),
\qquad \tilde{x}_{ij} > 0.
\]

\bigskip

Le passage au logarithme montre qu'il s'agit d'une somme pondérée sur une échelle logarithmique :
la substituabilité entre métriques devient limitée, et une valeur proche de zéro tire tout le
score vers le bas. Deux précautions : décaler les données ($\tilde{x} \mapsto \tilde{x} + \epsilon$
avec $\epsilon$ petit, par exemple $10^{-3}$) pour éviter les zéros, et vérifier que le résultat
n'est pas piloté par ce décalage --- c'est un test de sensibilité à part entière.

\bigskip

%-----------------------------------------------------------------------
\subsection{Indice à pénalité de déséquilibre (Mazziotta--Pareto)}
\label{sec:mpi}

\medskip

\paragraph{Idée.} Pénaliser explicitement l'hétérogénéité du profil d'un produit, sans introduire
de poids : à moyenne égale, un produit vulnérable sur une seule dimension et sûr ailleurs est
jugé plus sévèrement qu'un produit uniformément moyen.

\bigskip

\paragraph{Construction.} On normalise en scores standardisés recalés sur une moyenne de $100$ et
un écart-type de $10$:

\medskip

\[
z_{ij} = 100 + 10\,\frac{\tilde{x}_{ij} - \bar{x}_j}{\sigma_j}.
\]

\bigskip

Pour chaque produit, on note $M_i$ et $S_i$ la moyenne et l'écart-type \emph{en ligne} des
$(z_{ij})_j$, et $\mathrm{cv}_i = S_i/M_i$ le coefficient de variation horizontal. L'indice
s'écrit alors, dans sa version adaptée à une grandeur dont l'augmentation est défavorable:

\medskip

\[
\mathrm{MPI}^{+}_i = M_i + S_i \cdot \mathrm{cv}_i = M_i + \frac{S_i^2}{M_i}.
\]

\bigskip

Le terme $S_i^2/M_i$ est la pénalité: nulle pour un profil parfaitement équilibré, croissante
avec la dispersion des métriques du produit. Aucun coefficient n'est à fixer, ce qui en fait une
des rares agrégations scalaires entièrement dépourvue de paramètre libre.

\bigskip

\begin{avertissement}
  $\mathrm{MPI}^{+}$ \emph{n'est pas monotone} au sens de la propriété~\ref{prop:monotonie}:
  augmenter la plus faible des métriques d'un produit réduit $S_i$, et peut donc faire baisser
  l'indice. Cette violation est intrinsèque à toute pénalité de déséquilibre. Elle doit être
  mesurée (§\ref{sec:violation}) et assumée, ou la méthode écartée.
\end{avertissement}

\bigskip

%-----------------------------------------------------------------------
\subsection{TOPSIS}
\label{sec:topsis}

\paragraph{Idée.} Classer les produits par leur proximité relative à un pôle « idéal » et leur
éloignement d'un pôle « anti-idéal », tous deux construits empiriquement à partir des données.
Ici, le pôle idéal est le produit fictif maximalement vulnérable.

\paragraph{Construction.} Sur la matrice normalisée par colonne
$v_{ij} = w_j\, \tilde{x}_{ij} / \sqrt{\sum_k \tilde{x}_{kj}^2}$, on pose

\medskip

\[
A^{+}_j = \max_i v_{ij},
\qquad
A^{-}_j = \min_i v_{ij},
\]
\[
D^{\pm}_i = \Big(\sum_{j=1}^d (v_{ij}-A^{\pm}_j)^2\Big)^{1/2},
\qquad
s^{\mathrm{TOPSIS}}(x_i) = \frac{D^{-}_i}{D^{+}_i + D^{-}_i} \in [0,1].
\]

\bigskip

Associée à des poids issus de l'entropie ou de CRITIC, la procédure devient entièrement
empirique. Elle reste sensible aux points extrêmes, puisque $A^{\pm}$ sont des minima et maxima;
une variante robuste utilise les quantiles $\hat F_j^{-1}(0{,}99)$ et $\hat F_j^{-1}(0{,}01)$.

\bigskip

%-----------------------------------------------------------------------
\subsection{Distance de Mahalanobis à l'anti-idéal}
\label{sec:maha}

\medskip

\paragraph{Idée.} Si les métriques sont corrélées, une distance euclidienne compte plusieurs fois
la même information. La distance de Mahalanobis corrige cet effet en blanchissant l'espace:
elle mesure les écarts en unités de dispersion \emph{du nuage}, dans chaque direction.

\medskip

\paragraph{Construction.} Soit $\Sigma$ la matrice de covariance des métriques et
$x^{-} = \big(\hat F_1^{-1}(0{,}01),\dots,\hat F_d^{-1}(0{,}01)\big)$ le pôle anti-idéal
(produit fictif le moins vulnérable):
\[
s^{\mathrm{Mah}}(x_i) = \sqrt{(x_i - x^{-})^\top \Sigma^{-1} (x_i - x^{-})}.
\]

\bigskip

En pratique, estimer $\Sigma$ par un estimateur robuste (déterminant de covariance minimum, MCD,
ou estimateur à rétrécissement de Ledoit--Wolf si $d$ est grand devant $n$), sans quoi les
produits extrêmes contamineront leur propre référence.

\bigskip

\begin{remarque}[Une somme pondérée déguisée]
  Écrire $\Sigma^{-1} = \Sigma^{-1/2}\Sigma^{-1/2}$ montre que $s^{\mathrm{Mah}}$ est une norme
  euclidienne sur les données blanchies $u_i = \Sigma^{-1/2}(x_i - x^{-})$. La méthode revient donc
  à donner un poids égal à chaque \emph{direction décorrélée}, et non à chaque métrique~: les poids
  implicites sur les métriques d'origine sont lus dans $\Sigma^{-1/2}$. Ils sont entièrement
  estimés sur les données, ce qui satisfait la contrainte, mais ils peuvent être négatifs --- avec
  la même conséquence sur la monotonie que dans le cas de l'ACP. C'est un point à vérifier
  systématiquement.
\end{remarque}

\bigskip

Cette remarque est aussi la porte d'entrée de la section suivante : la distance de Mahalanobis
est le cas particulier « elliptique » d'une construction plus générale fondée sur le transport
optimal.

\bigskip

%=======================================================================
\section{Rangs de Kantorovitch et transport optimal}
\label{sec:ot}
%=======================================================================

\medskip

Cette section répond à la question posée sur l'article de Klein, Bethune, Ndiaye et
Cuturi~\cite{klein2025}. La réponse tient en trois temps~: le cadre des rangs \emph{center-outward}
est effectivement transposable, mais la partie « prédiction conforme » de l'article ne l'est pas
directement, et le rang tel qu'il est défini dans l'article mesure l'\emph{atypicité}, pas la
vulnérabilité --- il faut le réorienter.

%-----------------------------------------------------------------------
\subsection{Le cadre \emph{center-outward}}

\medskip

\paragraph{Idée.} En dimension $1$, la fonction de répartition $F$ est l'unique application
croissante qui transporte la loi $P$ sur la loi uniforme : c'est ce qui fait du rang une notion
canonique. En dimension $d$, la généralisation de la monotonie est le gradient d'une fonction
convexe, et le théorème de Brenier fournit exactement cet objet. On transporte donc la loi des
observations sur la loi uniforme de la \emph{boule unité}, choisie comme référence parce qu'elle
est centrée, symétrique et bornée.

\bigskip

\begin{definition}[Loi \emph{center-outward}, rang et signe]
  \label{def:co}
  Soit $Z \sim P$ à valeurs dans $\R^d$, $P$ absolument continue. La \emph{fonction de répartition
  center-outward} de $P$ est l'application $\Tmap = \nabla\varphi$, gradient d'une fonction convexe,
  qui pousse $P$ sur la loi uniforme $\Unif$ de la boule unité $\Bball(0,1)$. On définit
  \[
  \underbrace{\mathrm{R}(z) = \norm{\Tmap(z)} \in [0,1]}_{\text{rang}},
  \qquad
  \underbrace{\mathrm{S}(z) = \Tmap(z)/\norm{\Tmap(z)} \in \Sphere^{d-1}}_{\text{signe}}.
  \]
  Les régions quantiles sont les images réciproques des boules,
  $\mathcal{R}_\alpha = \{z : \norm{\Tmap(z)} \le 1-\alpha\}$, et vérifient
  $\PP(Z \in \mathcal{R}_\alpha) = 1-\alpha$.
\end{definition}

\bigskip

\paragraph{Ce que cela apporte réellement.} Deux propriétés, et elles sont fortes :

\medskip

\begin{itemize}[nosep,leftmargin=1.4em]
  \item \textbf{Absence de distribution.} Les rangs empiriques $\norm{\Tmap_n(z_i)}$ sont
  uniformément répartis sur la grille $\{1/n_R, 2/n_R,\dots,1\}$ \emph{quelle que soit} la loi
  $P$. C'est l'analogue multivarié exact de la transformation en rangs : échelles hétérogènes,
  asymétries, queues lourdes et corrélations sont neutralisées d'un seul coup, sans qu'aucune
  décision de normalisation n'ait à être prise coordonnée par coordonnée.
  \item \textbf{Prise en compte de la géométrie.} Contrairement à une normalisation marginale
  suivie d'une norme, la construction tient compte de la forme du nuage : une valeur élevée dans
  une direction où le nuage est étalé est jugée moins remarquable que la même valeur dans une
  direction où il est resserré.
\end{itemize}

\bigskip

%-----------------------------------------------------------------------
\subsection{L'objection décisive : le rang mesure l'atypicité}
\label{sec:objection}

\medskip

Le rang $\norm{\Tmap(z)}$ est \emph{non orienté}. Il croît du centre du nuage vers sa périphérie,
dans \emph{toutes} les directions à la fois. Un produit très vulnérable sur toutes les métriques
et un produit remarquablement peu vulnérable sur toutes les métriques sont tous deux à la
périphérie du nuage : ils reçoivent le même rang. Utilisé tel quel, $\norm{\Tmap(z)}$ classe donc
les produits par \emph{singularité}, ce qui ne coïncide avec la vulnérabilité que si l'on postule
« vulnérable $=$ atypique » --- un postulat rarement défendable pour une nomenclature douanière,
où l'on cherche des valeurs élevées dans une direction privilégiée.

\medskip

\begin{figure}[htbp]
\centering
\begin{tikzpicture}[scale=1.0]
  % panneau gauche : nuage et contours center-outward
  \begin{scope}
    \node[font=\small\bfseries] at (2.4,5.0) {Rang non orienté $\norm{\Tmap(z)}$};
    \draw[->,Gris] (-0.3,0) -- (4.7,0) node[right,black,font=\scriptsize] {$m_1$};
    \draw[->,Gris] (0,-0.3) -- (0,4.2) node[above,black,font=\scriptsize] {$m_2$};
    \foreach \r in {0.55,1.05,1.55} {
      \draw[Gris,rotate around={32:(2.2,2.0)}] (2.2,2.0) ellipse ({\r*1.5} and {\r*0.75});
    }
    \foreach \x/\y in {1.5/1.4, 2.6/2.4, 2.0/1.8, 2.9/2.6, 1.9/1.5, 2.5/2.1, 3.0/2.9,
                       1.6/1.7, 2.3/2.3, 2.8/2.2, 1.8/1.3, 2.4/2.6}
      \fill[Gris] (\x,\y) circle (1.7pt);
    \fill[Rouge] (3.9,3.2) circle (3pt) node[right,font=\scriptsize,black] {$a$};
    \fill[Rouge] (0.5,0.8) circle (3pt) node[left,font=\scriptsize,black] {$b$};
    \node[font=\scriptsize,align=center,Rouge] at (2.2,-0.85)
      {$\norm{\Tmap(a)} = \norm{\Tmap(b)}$ : \\ indiscernables};
  \end{scope}
  % panneau droit : score oriente
  \begin{scope}[xshift=7.2cm]
    \node[font=\small\bfseries] at (2.4,5.0) {Score orienté $\scal{\Tmap(z)}{u^\star}$};
    \draw[->,Gris] (-0.3,0) -- (4.7,0) node[right,black,font=\scriptsize] {$m_1$};
    \draw[->,Gris] (0,-0.3) -- (0,4.2) node[above,black,font=\scriptsize] {$m_2$};
    \foreach \x/\y in {1.5/1.4, 2.6/2.4, 2.0/1.8, 2.9/2.6, 1.9/1.5, 2.5/2.1, 3.0/2.9,
                       1.6/1.7, 2.3/2.3, 2.8/2.2, 1.8/1.3, 2.4/2.6}
      \fill[Gris] (\x,\y) circle (1.7pt);
    \draw[Vert,very thick,-{Stealth}] (2.2,2.0) -- (3.5,3.3) node[above left,font=\scriptsize,black] {$u^\star$};
    \foreach \c in {-0.9,0,0.9,1.8} {
      \draw[Bleu,dashed,shift={(\c*0.7071,\c*0.7071)}] (0.3,3.9) -- (3.9,0.3);
    }
    \fill[Rouge] (3.9,3.2) circle (3pt) node[right,font=\scriptsize,black] {$a$};
    \fill[Rouge] (0.5,0.8) circle (3pt) node[left,font=\scriptsize,black] {$b$};
    \node[font=\scriptsize,align=center,Vert] at (2.2,-0.85)
      {$\scal{\Tmap(a)}{u^\star} \gg \scal{\Tmap(b)}{u^\star}$ : \\ ordre restauré};
  \end{scope}
\end{tikzpicture}
\caption{À gauche, les régions quantiles \emph{center-outward} et deux produits diamétralement
opposés qui reçoivent le même rang. À droite, la projection sur la direction de vulnérabilité
$u^\star$ rétablit l'orientation.}
\label{fig:orientation}
\end{figure}

%-----------------------------------------------------------------------
\subsection{Score de Kantorovitch orienté}
\label{sec:otscore}

\medskip

Le remède est fourni par la définition \ref{def:co} elle-même : la construction de Hallin
\emph{et al.} produit un rang \emph{et} un signe, et l'article n'exploite que le rang parce qu'il
construit des régions de couverture symétriques. Un problème de classement doit au contraire
utiliser les deux.

\medskip

\paragraph{Étape 1 --- Déterminer la direction de vulnérabilité.} On construit un produit fictif
« maximalement vulnérable » $x^{+}$, par exemple $x^{+}_j = \hat F_j^{-1}(0{,}99)$ pour tout $j$,
et l'on définit
\[
u^\star = \frac{\bar{\Tmap}(x^{+})}{\norm{\bar{\Tmap}(x^{+})}} \in \Sphere^{d-1},
\]
où $\bar{\Tmap}$ est l'extension hors échantillon de la carte de transport (§\ref{sec:entropique}).
Cette direction est \emph{entièrement estimée sur les données}~: aucune hiérarchie entre critères
n'est introduite, seul le sens commun « plus vulnérable est plus vulnérable » est utilisé.

\medskip

\paragraph{Étape 2 --- Projeter.} Trois scores possibles, par ordre croissant de sévérité:

\medskip

\begin{align*}
  \text{projection signée} \qquad
  & s^{\mathrm{OT}}_{\mathrm{proj}}(x_i) = \scal{\Tmap(x_i)}{u^\star}, \\[0.3em]
  \text{rang restreint à un cône} \qquad
  & s^{\mathrm{OT}}_{\mathrm{cone}}(x_i) = \norm{\Tmap(x_i)}\cdot
    \ind{\scal{\mathrm{S}(x_i)}{u^\star} \ge \cos\theta_0}, \\[0.3em]
  \text{rang modulé par le signe} \qquad
  & s^{\mathrm{OT}}_{\mathrm{mult}}(x_i) = \norm{\Tmap(x_i)}\cdot
    \frac{1+\scal{\mathrm{S}(x_i)}{u^\star}}{2}.
\end{align*}

\bigskip

Le premier est le plus simple et se lit comme une somme pondérée dans l'espace transporté, à
coefficients $u^\star$ estimés. Le deuxième donne un ensemble d'alerte plutôt qu'un classement, en
ne retenant que les produits périphériques \emph{du bon côté}, avec un demi-angle $\theta_0$
paramétrable. Le troisième interpole continûment entre les deux.

\medskip

\begin{algorithm}[htbp]
  \caption{Score de Kantorovitch orienté}
  \label{algo:otscore}
  \begin{algorithmic}[1]
    \Require $\tilde{X}\in\R^{n\times d}$, régularisation $\varepsilon$, taille de grille $m$, quantile de pôle $q$
    \Ensure scores $s^{\mathrm{OT}} \in \R^n$
    \State \textbf{Grille sphérique.} Tirer $w_1,\dots,w_m$ suite à faible discrépance sur $[0,1]^d$ (Sobol)
    \State $\theta_\ell \gets \Phi^{-1}(w_\ell)/\norm{\Phi^{-1}(w_\ell)}_2$
          \Comment{Directions quasi uniformes sur $\Sphere^{d-1}$}
    \State Tirer les rayons $\rho_\ell$ par échantillonnage stratifié sur $[0,1]$
    \State $u_\ell \gets \rho_\ell^{1/d}\,\theta_\ell$
          \Comment{Loi uniforme sur la boule, et non sur la sphère}
    \State \textbf{Carte de transport.} Résoudre le problème de Sinkhorn entre $\{\tilde{x}_i\}_{i\in[n]}$ et $\{u_\ell\}_{\ell\in[m]}$
          à la régularisation $\varepsilon$, obtenir les potentiels duaux $f^\star, g^\star$
    \State $\bar{\Tmap}(z) \gets \sum_{\ell=1}^m p_\ell(z)\, u_\ell$ \textbf{avec}
          $p_\ell(z) \propto \exp\!\big(-(\norm{z-u_\ell}^2 - g^\star_\ell)/\varepsilon\big)$
          \Comment{Carte entropique, valide hors échantillon}
    \State \textbf{Direction de vulnérabilité.} $x^{+}_j \gets \hat F_j^{-1}(q)$ pour tout $j$~;\quad
          $u^\star \gets \bar{\Tmap}(x^{+})/\norm{\bar{\Tmap}(x^{+})}$
    \For{$i = 1$ \textbf{à} $n$}
      \State $s^{\mathrm{OT}}_i \gets \scal{\bar{\Tmap}(\tilde{x}_i)}{u^\star}$
    \EndFor
    \State \textbf{Contrôle.} Vérifier que $\{\norm{\bar{\Tmap}(\tilde{x}_i)}\}_i$ est approximativement
          uniforme sur $[0,1]$ (test de Kolmogorov--Smirnov)
    \State \Return $s^{\mathrm{OT}}$
  \end{algorithmic}
\end{algorithm}

%-----------------------------------------------------------------------
\subsection{Quand cette approche apporte-t-elle quelque chose ?}
\label{sec:otvsmaha}

La remarque 3.4 de l'article donne le critère de décision, et il est très opérationnel: si la
loi source et la loi cible sont toutes deux elliptiques, la carte de Brenier est \emph{affine},
et le score de transport optimal se réduit exactement à une distance de Mahalanobis.

\medskip

\begin{remarque}[Critère de décision empirique]
  Le transport optimal n'apporte rien de plus que le §\ref{sec:maha} lorsque le nuage des métriques
  est approximativement elliptique. Il n'apporte quelque chose que si ce nuage est nettement
  non elliptique : asymétrie marquée, multimodalité (par exemple des régimes distincts selon les
  chapitres de la nomenclature), dépendance de queue asymétrique entre HHI et indices de
  dépendance. \textbf{Ce point se teste avant de payer le coût d'implémentation} : test d'ellipticité
  de la loi jointe, examen des projections bivariées, comparaison de la corrélation de rangs entre
  $s^{\mathrm{OT}}_{\mathrm{proj}}$ et $s^{\mathrm{Mah}}$. Si le $\tau$ de Kendall entre les deux
  dépasse, disons, $0{,}95$, la complexité supplémentaire n'est pas justifiable.
\end{remarque}
% /!\ Je ne sais pas si je suis d'accord avec cette remarque, si la méthode la plus complexe est plus généralement juste, alors, même si elle est computationnellement plus couteuse, je souhaite la conserver.

\bigskip

%-----------------------------------------------------------------------
\subsection{Estimation : carte entropique et hyperparamètres}
\label{sec:entropique}

\medskip

Le calcul exact de $\Tmap_n$ passe par un problème d'affectation optimale, résolu par l'algorithme
hongrois en $O(n^3)$, et impose une cible de même taille que l'échantillon. L'article contourne
les deux difficultés par la \emph{carte entropique} de Pooladian et Niles-Weed. On résout le
problème de Sinkhorn régularisé :

\medskip

\[
f^\star, g^\star = \arg\max_{f \in \R^n, g \in \R^m}\;
\scal{f}{\tfrac{\mathbf{1}_n}{n}} + \scal{g}{\tfrac{\mathbf{1}_m}{m}}
- \varepsilon \scal{e^{f/\varepsilon}}{K e^{g/\varepsilon}},
\qquad K_{i\ell} = e^{-\norm{z_i - u_\ell}^2/\varepsilon},
\]

\bigskip

puis on définit l'estimateur, valide hors échantillon,
\[
\bar{\Tmap}_\varepsilon(z) = \sum_{\ell=1}^{m} p_\ell(z)\, u_\ell,
\qquad
p_\ell(z) = \frac{\exp\!\big(-(\norm{z - u_\ell}^2 - g^\star_\ell)/\varepsilon\big)}
{\sum_{k=1}^{m}\exp\!\big(-(\norm{z - u_k}^2 - g^\star_k)/\varepsilon\big)}.
\]

\bigskip

Le coût est $O(nm)$ à l'ajustement et $O(m)$ par point à transporter.

\paragraph{Choix des hyperparamètres.} Les résultats empiriques de l'article convergent :
$\varepsilon = 0{,}1$ avec les normalisations automatiques d'OTT-JAX est un réglage robuste, et
$m$ entre $2^{12}$ et $2^{15}$ points de grille. Les deux paramètres doivent être \emph{augmentés
avec la dimension}. Le compromis est clair : $\varepsilon$ trop petit reproduit les mauvaises
propriétés statistiques du transport non régularisé, $\varepsilon$ trop grand fait s'effondrer la
carte estimée vers une image floue de la sphère. Le débiaisage des sorties de Sinkhorn n'améliore
pas les résultats.

\bigskip

\paragraph{Grille sphérique.} La qualité de la cible compte. On suit l'approche gaussienne : à
partir d'une suite à faible discrépance $w_1,\dots,w_m$ sur $[0,1]^d$, on pose
$\theta_\ell = \Phi^{-1}(w_\ell)/\norm{\Phi^{-1}(w_\ell)}_2$, qui fournit des directions quasi
uniformes sur $\Sphere^{d-1}$, combinées à des rayons stratifiés.

\bigskip

\begin{avertissement}[Malédiction de la dimension]
  L'estimation de cartes de transport se dégrade rapidement avec $d$. Les résultats de l'article
  montrent un avantage net jusqu'à $d \le 6$ et un avantage qui s'évanouit pour $14 \le d \le 16$.
  Un dispositif à trois ou quatre métriques (HHI, indices de dépendance) est donc exactement dans
  le régime favorable ; au-delà de six ou sept métriques, il faudra réduire la dimension au
  préalable ou renoncer.
\end{avertissement}

\bigskip

%-----------------------------------------------------------------------
\subsection{Ce qui n'est pas transposable, et ce qui l'est quand même}

\medskip

\paragraph{La prédiction conforme, telle quelle, ne s'applique pas.} L'article construit des
ensembles de prédiction pour la réponse $y_{n+1}$ d'un modèle $\hat y$, à partir de scores de
non-conformité $S(x,y,\hat y)$ --- typiquement des résidus. La garantie de couverture porte sur
une observation \emph{future}, sous hypothèse d'échangeabilité. Or le problème posé ici est
descriptif~: on dispose de la population complète des positions NC et on souhaite les ordonner.
Il n'y a ni modèle prédictif, ni observation future, ni résidus. La partie « conformalisation »
(propositions 3.2 et 3.6 de l'article) est donc hors sujet pour le classement lui-même.

\paragraph{Elle redevient pertinente pour calibrer un seuil.} Il existe toutefois un usage
légitime, et il est intéressant. Si l'on veut définir un \emph{seuil d'alerte} plutôt qu'un
palmarès --- « quels produits sont anormalement vulnérables ? » --- alors la région quantile
calibrée fournit exactement cela, avec une garantie valable à distance finie et sans hypothèse de
loi :

\medskip

\[
\widehat{\mathcal{R}}_\alpha = \big\{z : \norm{\bar{\Tmap}(z)} \le \hat r_\alpha \big\},
\qquad
\hat r_\alpha = \inf\big\{ r \ge 0 : \widehat{\Unif}(\Bball(0,r)) \ge 1 - \alpha \big\},
\]

\bigskip

l'ensemble d'alerte étant le complémentaire de $\widehat{\mathcal{R}}_\alpha$ \emph{intersecté}
avec le demi-espace $\scal{\mathrm{S}(z)}{u^\star} > 0$. L'intérêt est réel dans un contexte
administratif : le seuil n'est plus un « top~50 » arbitraire, mais un niveau $\alpha$ explicite,
avec une couverture garantie. Les usages où cette garantie a une vraie valeur :

\medskip

\begin{itemize}[nosep,leftmargin=1.4em]
  \item \textbf{Suivi temporel.} La carte est estimée sur l'année $t$, et l'on classe les produits
  de l'année $t+1$~: la garantie de couverture s'applique alors, l'échangeabilité étant l'hypothèse
  explicite (et testable) que la structure de vulnérabilité n'a pas changé. Un dépassement
  significatif du taux d'alerte attendu devient un signal de rupture de régime.
  \item \textbf{Nouvelles positions NC.} Une position créée en cours d'année reçoit un score et un
  statut d'alerte calibrés sur l'existant, sans réestimation.
\end{itemize}

\bigskip

%=======================================================================
\section{Comparer les agrégations et contrôler leur cohérence}
\label{sec:comparaison}
%=======================================================================

\medskip

Chaque méthode des sections précédentes produit un classement. Rien ne garantit qu'ils
concordent, et c'est justement ce désaccord qui porte l'information la plus utile : il délimite
la part du résultat qui tient aux données et la part qui tient au choix méthodologique. Cette
section propose un protocole en cinq questions.

\medskip

\begin{figure}[htbp]
  \centering
  \begin{tikzpicture}[
    node distance=0.6cm and 1.0cm,
    bloc/.style={draw=Bleu,fill=Bleu!6,rounded corners=2pt,align=center,font=\small,
                minimum height=1.0cm,text width=3.4cm,inner sep=4pt},
    outil/.style={font=\scriptsize,Gris,text width=3.0cm,align=right},
    fl/.style={-{Stealth[length=2.5mm]},Gris,thick}
  ]
    \node[bloc] (q1) {\textbf{Q1} --- Les classements\\ sont-ils d'accord ?};
    \node[bloc,below=of q1] (q2) {\textbf{Q2} --- D'accord\\ \emph{en tête} ?};
    \node[bloc,below=of q2] (q3) {\textbf{Q3} --- Cohérents avec\\ la dominance ?};
    \node[bloc,below=of q3] (q4) {\textbf{Q4} --- Stables au\\ rééchantillonnage ?};
    \node[bloc,below=of q4] (q5) {\textbf{Q5} --- Non réductibles\\ à une seule métrique ?};
    \foreach \a/\b in {q1/q2,q2/q3,q3/q4,q4/q5} \draw[fl] (\a) -- (\b);

    \node[outil,left=0.35cm of q1] {$\tau_b$ et $W$ de Kendall};
    \node[outil,left=0.35cm of q2] {RBO, Jaccard sur le top-$k$};
    \node[outil,left=0.35cm of q3] {taux de violation $V(s)$};
    \node[outil,left=0.35cm of q4] {bootstrap, SMAA};
    \node[outil,left=0.35cm of q5] {\emph{leave-one-metric-out}};

    \node[bloc,fill=Vert!8,draw=Vert,right=1.9cm of q3,text width=3.8cm]
        (res) {\textbf{Synthèse}\\ classement consensus\\ $+$ intervalles de rang\\ $+$ cas litigieux};
    \draw[fl] (q1.east) to[out=0,in=90] (res.north);
    \draw[fl] (q3.east) -- (res.west);
    \draw[fl] (q5.east) to[out=0,in=270] (res.south);
  \end{tikzpicture}
  \caption{Protocole de contrôle de cohérence.}
\end{figure}

\bigskip

%-----------------------------------------------------------------------
\subsection{Q1 --- Concordance globale des classements}

\medskip

Soit $r^{(A)}, r^{(B)} \in [n]^n$ les vecteurs de rangs produits par deux méthodes $A$ et $B$.

\paragraph{Corrélations de rangs.} Le $\tau_b$ de Kendall est préférable au $\rho$ de Spearman
ici : il s'interprète directement en proportion de paires concordantes, gère les ex \ae quo et est
moins sensible aux écarts sur les rangs médians, dont on se soucie peu. On produit la matrice
$T = (\tau_b(r^{(A)}, r^{(B)}))_{A,B}$ pour toutes les paires de méthodes, visualisée en carte de
chaleur et complétée par une classification ascendante hiérarchique sur $1-\tau_b$ : les groupes
de méthodes ainsi révélés correspondent en général aux grandes options méthodologiques
(compensatoire ou non, décorrélée ou non).

\bigskip

\paragraph{Concordance d'ensemble.} Le coefficient $W$ de Kendall mesure l'accord simultané de
$K$ classements :
\[
W = \frac{12 \sum_{i=1}^n \big(R^{\Sigma}_i - \tfrac{K(n+1)}{2}\big)^2}{K^2 n (n^2-1)}
\;\in [0,1],
\qquad R^{\Sigma}_i = \sum_{k=1}^K r^{(k)}_i,
\]

\bigskip

avec $W = 1$ pour un accord parfait. Un $W$ élevé signifie que la question du choix de méthode est
secondaire ; un $W$ faible signifie qu'aucun classement unique ne devrait être publié sans
mention de son incertitude.

\bigskip

\paragraph{Déplacement moyen de rang.} La mesure la plus lisible pour un lecteur non
statisticien, recommandée par le \emph{Handbook} de l'OCDE :

\medskip

\[
\bar{R}_{A\to B} = \frac{1}{n}\sum_{i=1}^n \abs{r^{(A)}_i - r^{(B)}_i},
\]
« en changeant de méthode, un produit se déplace en moyenne de tant de places ».

\bigskip

%-----------------------------------------------------------------------
\subsection{Q2 --- Concordance en tête de classement}

\medskip

Le $\tau$ global est trompeur pour un usage opérationnel : seuls les premiers rangs seront
exploités, et deux classements peuvent avoir un $\tau$ de $0{,}9$ tout en partageant la moitié
seulement de leur top 50.

\medskip

\paragraph{Recouvrement du top-$k$.} $J_k(A,B) = \abs{\mathcal{T}^{A}_k \cap \mathcal{T}^{B}_k}/k$,
tracé en fonction de $k$ : la courbe montre à partir de quelle profondeur les méthodes commencent
à s'accorder.

\medskip

\paragraph{\emph{Rank-biased overlap}.} Le RBO agrège ces recouvrements en pondérant
géométriquement les profondeurs, ce qui donne beaucoup plus de poids aux premiers rangs :
\[
\mathrm{RBO}(A,B\,;\,p) = (1-p)\sum_{k=1}^{\infty} p^{\,k-1}\, J_k(A,B),
\qquad p \in (0,1).
\]

\bigskip

Le paramètre $p$ se règle par la profondeur d'intérêt : l'espérance du rang examiné vaut
$1/(1-p)$, donc $p = 0{,}98$ pour une attention centrée sur les cinquante premiers.

\bigskip

%-----------------------------------------------------------------------
\subsection{Q3 --- Cohérence avec la dominance : le test dur}
\label{sec:violation}

\medskip

C'est le seul contrôle qui ne repose sur aucune convention et qui puisse invalider une méthode.

\medskip

\begin{definition}[Taux de violation de la dominance]
  Soit $\mathcal{D} = \{(i,k) : x_i \pareto x_k\}$ l'ensemble des paires comparables au sens de
  Pareto. Le taux de violation d'un score $s$ est :

  \medskip

  \[
  V(s) = \frac{1}{\abs{\mathcal{D}}}
  \sum_{(i,k)\in\mathcal{D}} \ind{s(x_i) \le s(x_k)}.
  \]
\end{definition}

\bigskip

Un score $V(s) > 0$ inverse des paires sur lesquelles toutes les métriques sont d'accord : c'est
indéfendable devant un comité, quelle que soit la sophistication de la méthode. On attend :

\medskip

\begin{itemize}
  \item $V = 0$ pour le comptage de dominance, la somme pondérée à poids positifs, la moyenne
  géométrique, TOPSIS, BoD ;
  \item $V > 0$ possible pour l'ACP et Mahalanobis dès qu'apparaît une saturation négative ;
  \item $V > 0$ structurellement pour $\mathrm{MPI}^{+}$ (pénalité de déséquilibre) et pour le
  score de transport optimal, dont la carte est monotone \emph{cycliquement} mais non
  coordonnée par coordonnée.
\end{itemize}

\bigskip

Il faut donc rapporter $V(s)$ pour chaque méthode et, pour celles qui violent, examiner les paires
concernées : si les violations se concentrent sur des paires quasi indifférentes (écarts
inférieurs à l'incertitude d'estimation des métriques), elles sont sans portée pratique~; si elles
touchent des paires nettement séparées, la méthode doit être écartée.

\bigskip

\paragraph{Contrôle associé.} Vérifier que le front $\mathcal{F}_1$ est bien concentré dans le
haut de chaque classement~: rapporter le rang médian et le rang maximal des produits non dominés.

\medskip

\begin{algorithm}[htbp]
  \caption{Tableau de bord de cohérence}
  \label{algo:coherence}
  \begin{algorithmic}[1]
    \Require $\tilde{X}$, ensemble de scores $\{s^{(1)},\dots,s^{(K)}\}$, profondeur $k$, paramètre RBO $p$
    \Ensure matrices de concordance, taux de violation, classement consensus
    \State $\mathcal{D} \gets$ paires en relation de dominance (algorithme~\ref{algo:nds})
    \For{$a = 1$ \textbf{à} $K$}
      \State $r^{(a)} \gets \operatorname{rang}(-s^{(a)})$ \Comment{Rang $1$ au plus vulnérable}
      \State $V_a \gets \abs{\mathcal{D}}^{-1}\sum_{(i,k)\in\mathcal{D}} \ind{s^{(a)}_i \le s^{(a)}_k}$
      \For{$b = 1$ \textbf{à} $K$}
        \State $T_{ab} \gets \tau_b(r^{(a)}, r^{(b)})$~;\quad
              $O_{ab} \gets \mathrm{RBO}(r^{(a)}, r^{(b)}; p)$
      \EndFor
    \EndFor
    \State $W \gets$ coefficient de concordance de Kendall sur $\{r^{(a)}\}_a$
    \State $\bar{R}_i \gets K^{-1}\sum_a r^{(a)}_i$ \Comment{Rang de Borda, classement consensus}
    \State $\mathcal{L} \gets \{ i : \max_a r^{(a)}_i - \min_a r^{(a)}_i > \eta \}$
          \Comment{Produits litigieux, à examiner individuellement}
    \State \Return $T$, $O$, $V$, $W$, $\bar{R}$, $\mathcal{L}$
  \end{algorithmic}
\end{algorithm}

%-----------------------------------------------------------------------
\subsection{Q4 --- Stabilité}
\label{sec:bootstrap}

\medskip

\paragraph{Bootstrap sur les produits.} Rééchantillonner $B$ fois (typiquement $B = 1\,000$) les
$n$ produits avec remise, recalculer intégralement la chaîne --- \emph{y compris les poids}, dont
c'est tout l'enjeu puisqu'ils sont estimés --- et collecter la distribution du rang de chaque
produit. On rapporte un intervalle de rang à $90\,\%$ plutôt qu'un rang ponctuel. Un produit dont
l'intervalle est $[3, 412]$ ne doit pas être présenté comme « troisième plus vulnérable ».
% /!\ Comment introduire un intervalle de confiance autour du rang ?

\medskip

\paragraph{Propagation de l'incertitude des métriques.} Si les HHI et indices de dépendance sont
eux-mêmes estimés sur un nombre fini de déclarations, leur erreur d'échantillonnage se propage.
Simuler $x_{ij}^{(b)} \sim \mathcal{N}(\hat{x}_{ij}, \hat{\sigma}^2_{ij})$, ou rééchantillonner
directement les déclarations sous-jacentes, et recalculer le classement. C'est souvent la source
d'instabilité dominante pour les positions à faibles volumes.

\medskip

\paragraph{Stabilité aux choix méthodologiques.} Faire varier systématiquement le schéma de
normalisation ($\times 5$), la méthode de pondération ($\times 4$) et la fonction d'agrégation
($\times 5$), soit une centaine de variantes. Pour chaque produit, tracer la distribution des
rangs obtenus. On peut aller plus loin et décomposer la variance du rang par indices de Sobol
attribués à chacun des trois facteurs : le résultat --- « $62\,\%$ de la variance du rang vient du
choix de normalisation » --- est précieux, car il indique où concentrer l'effort de justification.

\bigskip

%-----------------------------------------------------------------------
\subsection{SMAA : formaliser l'ignorance sur les poids}
\label{sec:smaa}

\medskip

\paragraph{Idée.} Plutôt que de choisir un vecteur de poids, on explore \emph{tout} l'espace des
poids admissibles et l'on rapporte, pour chaque produit, la probabilité d'occuper chaque rang.
C'est la formalisation exacte de la situation décrite : incapacité à hiérarchiser les critères,
donc loi uniforme sur le simplexe.

\medskip

\paragraph{Construction.} On tire $w^{(t)} \sim \Unif(\simplex)$, ce qui s'obtient exactement par
une loi de Dirichlet $\mathrm{Dir}(1,\dots,1)$, on calcule le classement induit, et l'on estime:

\medskip

\begin{align*}
  \text{indice d'acceptabilité de rang} \quad & b^{\,r}_i = \PP\big(\text{rang}(i) = r\big), \\
  \text{vecteur de poids central} \quad & w^{c}_i = \EE\big[w \mid \text{rang}(i) = 1\big], \\
  \text{facteur de confiance} \quad & p_i = \PP\big(\text{rang}(i) \le k \mid w \sim \Unif(\simplex)\big).
\end{align*}

\bigskip

La sortie utile pour un dossier administratif est $p_i$ avec $k$ fixé par la capacité d'action
(par exemple $k=50$) : « ce produit figure dans les cinquante plus vulnérables sous $94\,\%$ des
pondérations possibles ». Une telle affirmation est nettement plus solide, et plus honnête,
qu'un rang ponctuel.

\paragraph{Restrictions partielles.} SMAA accepte des contraintes d'ordre partiel sur les poids
($w_1 \ge w_2$, ou $w_j \ge 0{,}05$) que l'on peut ajouter si un consensus minimal existe, en
tirant uniformément dans le polytope ainsi défini par échantillonnage par rejet ou \emph{hit-and-run}.
On peut alors mesurer combien d'information une hiérarchisation même faible ferait gagner.

\medskip

\begin{algorithm}[htbp]
  \caption{SMAA --- indices d'acceptabilité de rang}
  \label{algo:smaa}
  \begin{algorithmic}[1]
    \Require $\tilde{X}\in\R^{n\times d}$, nombre de tirages $T$, fonction d'agrégation $g$, seuil $k$
    \Ensure indices $b \in [0,1]^{n\times n}$, poids centraux $w^c$, facteurs de confiance $p$
    \State $b \gets 0$~;\quad $c \gets 0 \in \R^{n\times d}$~;\quad $\nu \gets 0 \in \N^n$
    \For{$t = 1$ \textbf{à} $T$}
      \State $w^{(t)} \sim \mathrm{Dir}(1,\dots,1)$ \Comment{Loi uniforme sur le simplexe}
      \State $s^{(t)}_i \gets g(\tilde{x}_i ; w^{(t)})$ pour tout $i$
      \State $r^{(t)} \gets \operatorname{rang}(-s^{(t)})$
      \For{$i = 1$ \textbf{à} $n$}
        \State $b_{i, r^{(t)}_i} \gets b_{i, r^{(t)}_i} + 1/T$
        \If{$r^{(t)}_i = 1$} $c_i \gets c_i + w^{(t)}$~;\quad $\nu_i \gets \nu_i + 1$
        \EndIf
      \EndFor
    \EndFor
    \State $w^{c}_i \gets c_i / \max(\nu_i, 1)$~;\quad $p_i \gets \sum_{r \le k} b_{i r}$
    \State \Return $b$, $w^c$, $p$
  \end{algorithmic}
\end{algorithm}

\bigskip

\begin{remarque}
  Pour $n \sim 5\,000$ et $T \sim 10^4$, la matrice complète $b$ est de taille $5\times10^7$: ne
  conserver que les rangs $r \le k$ et l'agrégat $p$. L'ensemble se vectorise entièrement en NumPy
  (un produit matriciel $T \times d$ par $d \times n$, puis un \texttt{argsort} par ligne).
\end{remarque}

\bigskip

%-----------------------------------------------------------------------
\subsection{Q5 --- Le score est-il réductible à une seule métrique ?}

\medskip

Un risque courant des pondérations endogènes est de produire un score qui reproduit à peu de
choses près l'une des métriques d'entrée --- auquel cas tout l'appareil méthodologique n'a servi
à rien.

\medskip

\paragraph{Corrélations score--métrique.} Calculer $\tau_b(s, \tilde{x}_{\cdot j})$ pour tout $j$.
Si l'une dépasse nettement les autres, examiner pourquoi (poids dominant, métrique très
dispersée, corrélation avec toutes les autres).

\medskip

\paragraph{\emph{Leave-one-metric-out}.} Recalculer la chaîne complète en retirant chaque métrique
à tour de rôle, et mesurer $\tau_b$ et le recouvrement du top-$k$ avec le classement complet. Une
métrique dont le retrait ne change rien est soit redondante, soit neutralisée par la pondération~;
une métrique dont le retrait bouleverse tout porte à elle seule le résultat. Les deux situations
appellent un commentaire explicite dans la note de méthode.

\medskip

\paragraph{Décomposition des contributions.} Pour une agrégation linéaire, la contribution de la
métrique $j$ au score du produit $i$ est $w_j \tilde{x}_{ij} / s_i$~; on peut en dresser un profil
pour les produits de tête, ce qui rend le classement lisible et défendable individuellement.

\bigskip

%-----------------------------------------------------------------------
\subsection{Classement consensus}

Lorsqu'on doit tout de même publier un classement unique, trois constructions, toutes
indépendantes des poids :

\medskip

\begin{itemize}[leftmargin=1.4em,itemsep=2pt]
  \item \textbf{Borda} : $\bar{R}_i = \frac{1}{K}\sum_a r^{(a)}_i$. Immédiat, robuste, mais
  sensible aux méthodes fortement corrélées entre elles, qui votent alors plusieurs fois pour la
  même position. Corriger en pondérant chaque méthode par $1/(\text{taille de son groupe})$ dans
  la classification de Q1.
  \item \textbf{Copeland} : pour chaque paire $(i,k)$, compter le nombre de méthodes plaçant $i$
  devant $k$~; le score est le nombre de duels gagnés. Plus robuste aux valeurs de rang extrêmes.
  \item \textbf{Médiane de Kemeny} : le classement minimisant la somme des distances de Kendall aux
  $K$ classements. C'est le consensus le mieux fondé axiomatiquement, mais le problème est
  NP-difficile ; à $n \sim 5\,000$, on ne l'utilise que sur une présélection (les cent premiers du
  classement de Borda), résolue par programmation linéaire en nombres entiers.
\end{itemize}

\bigskip

\paragraph{Ce qu'il faut publier.} Un classement consensus, assorti pour chaque produit d'un
intervalle de rang bootstrap et de son facteur de confiance SMAA, plus la liste explicite des
produits litigieux $\mathcal{L}$ dont le rang dépend fortement de la méthode. Cette dernière liste
n'est pas un aveu de faiblesse~: c'est le résultat le plus informatif du dispositif, puisqu'elle
identifie exactement les produits pour lesquels une expertise humaine est nécessaire.

\bigskip

%-----------------------------------------------------------------------
\subsection{Validation externe}

Si l'on dispose d'un critère externe --- ruptures d'approvisionnement observées, alertes passées,
mesures de sauvegarde effectivement déclenchées, épisodes de tension documentés --- la validation
change de nature et devient une question de pouvoir prédictif~: AUC de la courbe ROC, précision
au rang $k$, gain cumulé actualisé. C'est le seul argument capable de départager objectivement
deux agrégations, et il vaut la peine de constituer un tel jeu d'événements même modeste
(quelques dizaines de cas avérés suffisent à discriminer). En son absence, il faut se résoudre à
ne conclure que sur la robustesse interne.

%=======================================================================
\section{Démarche recommandée}
\label{sec:demarche}
%=======================================================================

\begin{figure}[htbp]
\centering
\begin{tikzpicture}[
  node distance=0.75cm,
  etape/.style={draw=Bleu,fill=Bleu!7,rounded corners=3pt,align=center,font=\small,
                text width=10.5cm,minimum height=0.9cm,inner sep=5pt},
  sortie/.style={draw=Vert,fill=Vert!8,rounded corners=3pt,align=center,font=\small,
                 text width=10.5cm,minimum height=0.9cm,inner sep=5pt},
  fl/.style={-{Stealth[length=3mm]},Gris,very thick}
]
  \node[etape] (e1) {\textbf{1. Prétraitement} --- orientation des polarités, winsorisation,
        filtrage des positions à faible effectif, diagnostic de corrélation};
  \node[etape,below=of e1] (e2) {\textbf{2. Invariant sans hypothèse} --- front de Pareto
        $\mathcal{F}_1$ et comptage de dominance $\delta$};
  \node[etape,below=of e2] (e3) {\textbf{3. Scores concurrents} --- grille
        (normalisation $\times$ pondération endogène $\times$ agrégation), plus le score de
        Kantorovitch orienté si le nuage est non elliptique};
  \node[etape,below=of e3] (e4) {\textbf{4. Contrôles de cohérence} --- taux de violation de la
        dominance, $\tau_b$ et RBO entre méthodes, \emph{leave-one-metric-out}};
  \node[etape,below=of e4] (e5) {\textbf{5. Incertitude} --- bootstrap sur les produits et sur les
        métriques, SMAA sur le simplexe des poids};
  \node[sortie,below=of e5] (e6) {\textbf{Livrable} --- classement consensus, intervalles de rang,
        facteurs de confiance, liste des produits litigieux};
  \foreach \a/\b in {e1/e2,e2/e3,e3/e4,e4/e5,e5/e6} \draw[fl] (\a) -- (\b);
\end{tikzpicture}
\caption{Enchaînement recommandé. Les étapes 2 et 5 sont celles qui donnent au dispositif sa
défendabilité~; l'étape 3, prise isolément, n'en a aucune.}
\end{figure}

\paragraph{Ordre de priorité, en pratique.} Si le temps disponible est limité, l'ordre de valeur
marginale décroissante est le suivant~:
\begin{enumerate}[leftmargin=1.6em,itemsep=2pt]
  \item \textbf{Front de Pareto et comptage de dominance.} Coût nul, aucune hypothèse, résultat
  immédiatement publiable.
  \item \textbf{SMAA sur une somme pondérée.} Répond directement à l'impossibilité de hiérarchiser
  les critères, et fournit des énoncés probabilisés plutôt qu'un classement fragile.
  \item \textbf{Deux ou trois scores contrastés} --- par exemple somme pondérée par CRITIC,
  moyenne géométrique, et TOPSIS --- avec le tableau de bord de cohérence de
  l'algorithme~\ref{algo:coherence}.
  \item \textbf{Bootstrap} sur les métriques, en priorité pour les positions à faible volume.
  \item \textbf{Score de Kantorovitch orienté}, uniquement si le test d'ellipticité du
  §\ref{sec:otvsmaha} conclut à une géométrie franchement non gaussienne, et si $d \le 6$.
\end{enumerate}

\paragraph{Sur le transport optimal, en résumé.} L'approche est légitime et techniquement
transposable~; elle apporte une normalisation multivariée sans hypothèse de loi, ce qui est un
argument méthodologique fort. Mais elle exige une réorientation (§\ref{sec:otscore}) sans laquelle
elle mesure autre chose que ce qui est cherché, elle se réduit à une distance de Mahalanobis dès
que le nuage est elliptique, et elle introduit deux hyperparamètres à régler là où les méthodes
concurrentes n'en ont aucun. Elle mérite d'être testée, mais comme candidate parmi d'autres dans
le protocole de comparaison --- et non comme méthode de référence choisie \emph{a priori} pour son
élégance.

%=======================================================================
\section{Repères d'implémentation}
%=======================================================================

\begin{table}[htbp]
\centering
\small
\renewcommand{\arraystretch}{1.25}
\begin{tabularx}{\textwidth}{@{}l l X@{}}
\toprule
\textbf{Étape} & \textbf{Bibliothèque} & \textbf{Points d'attention} \\
\midrule
Normalisation & \texttt{scikit-learn} &
\texttt{QuantileTransformer} en sortie gaussienne pour la normalisation par quantile~; \texttt{RobustScaler} pour la variante médiane/MAD. \\
Front de Pareto & \texttt{paretoset}, \texttt{pymoo} &
\texttt{pymoo.util.nds.fast\_non\_dominated\_sort} pour le tri par couches. Version vectorisée
NumPy suffisante jusqu'à $n \sim 10^4$. \\
MCDA & \texttt{pymcdm}, \texttt{scikit-criteria} &
\texttt{pymcdm} implémente TOPSIS, PROMETHEE, ainsi que les pondérations par entropie et CRITIC. \\
ACP robuste & \texttt{scikit-learn}, \texttt{statsmodels} &
\texttt{MinCovDet} pour la covariance robuste~; \texttt{FactorAnalysis} $+$ rotation varimax pour
la procédure OCDE--JRC. \\
BoD / DEA & \texttt{scipy.optimize.linprog}, \texttt{PuLP} &
Réutiliser une même fabrique de contraintes pour les $n$ programmes~; paralléliser par
\texttt{joblib}. \\
Transport optimal & \texttt{ott-jax}, \texttt{POT} &
\texttt{ott-jax} offre directement les potentiels duaux et la carte entropique hors échantillon,
et se compile en JIT sur la grille sphérique. \\
Comparaison & \texttt{scipy.stats}, \texttt{rbo} &
\texttt{kendalltau}, \texttt{spearmanr}, \texttt{friedmanchisquare}. Le $W$ de Kendall se déduit
du $\chi^2$ de Friedman. \\
\bottomrule
\end{tabularx}
\caption{Écosystème Python.}
\end{table}

\noindent
Le score de Kantorovitch orienté est la seule brique sans implémentation directement disponible.
Le squelette suivant en donne la structure~:

\begin{lstlisting}
import jax.numpy as jnp
from jax.scipy.stats import norm
from ott.geometry import pointcloud
from ott.problems.linear import linear_problem
from ott.solvers.linear import sinkhorn
from scipy.stats import qmc


def spherical_uniform_grid(n_points, dim, seed=0):
    """Draw a low-discrepancy sample from the uniform measure on the unit ball.

    Args:
        n_points: Number of target points in the reference measure.
        dim: Ambient dimension of the score space.
        seed: Seed of the Sobol engine.

    Returns:
        Array of shape ``(n_points, dim)`` holding the reference grid.
    """
    # Suite de Sobol sur l'hypercube, puis transformation gaussienne
    engine = qmc.Sobol(d=dim + 1, scramble=True, seed=seed)
    sample = jnp.asarray(engine.random(n_points))
    sample = jnp.clip(sample, 1e-6, 1.0 - 1e-6)

    # Directions quasi uniformes sur la sphere unite
    gaussians = norm.ppf(sample[:, :dim])
    directions = gaussians / jnp.linalg.norm(gaussians, axis=1, keepdims=True)

    # Rayons stratifies : la puissance 1/dim assure l'uniformite dans la boule
    radii = sample[:, dim:] ** (1.0 / dim)
    return directions * radii


def oriented_kantorovich_score(x, epsilon=0.1, n_target=2**14, pole_quantile=0.99):
    """Rank products by an orientation-corrected center-outward score.

    The empirical distribution of the metrics is transported onto the uniform
    measure on the unit ball. Ranks alone are not orientable, so the transported
    points are projected on the image of a synthetic "maximally vulnerable"
    product, which restores a meaningful direction.

    Args:
        x: Array of shape ``(n, d)`` with all metrics in positive polarity.
        epsilon: Entropic regularisation strength.
        n_target: Size of the reference grid on the unit ball.
        pole_quantile: Marginal quantile defining the synthetic pole.

    Returns:
        Array of shape ``(n,)`` where a higher value means more vulnerable.
    """
    x = jnp.asarray(x)
    grid = spherical_uniform_grid(n_target, x.shape[1])

    # Resolution du probleme de Sinkhorn entre les scores et la boule uniforme
    geom = pointcloud.PointCloud(x, grid, epsilon=epsilon)
    solution = sinkhorn.Sinkhorn()(linear_problem.LinearProblem(geom))
    if not solution.converged:
        raise RuntimeError("Sinkhorn n'a pas converge : augmenter epsilon.")

    # Carte entropique, valide hors echantillon
    transport = solution.to_dual_potentials().transport

    # Direction de vulnerabilite, estimee sur les donnees
    pole = jnp.quantile(x, pole_quantile, axis=0)
    pole_image = transport(pole[None, :])[0]
    direction = pole_image / jnp.linalg.norm(pole_image)

    # Projection signee : rang et signe combines
    return transport(x) @ direction
\end{lstlisting}

\noindent
Deux contrôles à ajouter systématiquement~: la vérification que les normes
$\norm{\bar{\Tmap}(x_i)}$ suivent bien une loi uniforme sur $[0,1]$ (sans quoi la carte n'a pas
convergé vers un transport valide), et la comparaison du classement obtenu avec celui de la
distance de Mahalanobis, qui indique si le surcoût est justifié.

%=======================================================================
\begin{thebibliography}{99}
\small

\bibitem{klein2025}
M.~Klein, L.~Bethune, E.~Ndiaye, M.~Cuturi.
\emph{Multivariate Conformal Prediction using Optimal Transport}.
arXiv:2502.03609, 2025.

\bibitem{hallin2021}
M.~Hallin, E.~del Barrio, J.~Cuesta-Albertos, C.~Matrán.
\emph{Distribution and quantile functions, ranks and signs in dimension $d$: a measure
transportation approach}.
The Annals of Statistics, 49(2), 1139--1165, 2021.

\bibitem{chernozhukov2017}
V.~Chernozhukov, A.~Galichon, M.~Hallin, M.~Henry.
\emph{Monge--Kantorovich depth, quantiles, ranks and signs}.
The Annals of Statistics, 45(1), 223--256, 2017.

\bibitem{cuturi2013}
M.~Cuturi.
\emph{Sinkhorn distances: lightspeed computation of optimal transport}.
NeurIPS, 2292--2300, 2013.

\bibitem{pooladian2021}
A.-A.~Pooladian, J.~Niles-Weed.
\emph{Entropic estimation of optimal transport maps}.
arXiv:2109.12004, 2021.

\bibitem{peyre2019}
G.~Peyré, M.~Cuturi.
\emph{Computational Optimal Transport}.
Foundations and Trends in Machine Learning, 11(5--6), 2019.

\bibitem{thurin2025}
G.~Thurin, K.~Nadjahi, C.~Boyer.
\emph{Optimal transport-based conformal prediction}.
arXiv:2501.18991, 2025.

\bibitem{nardo2008}
M.~Nardo, M.~Saisana, A.~Saltelli, S.~Tarantola, A.~Hoffman, E.~Giovannini.
\emph{Handbook on Constructing Composite Indicators: Methodology and User Guide}.
OCDE et Centre commun de recherche de la Commission européenne, 2008.

\bibitem{saisana2005}
M.~Saisana, A.~Saltelli, S.~Tarantola.
\emph{Uncertainty and sensitivity analysis techniques as tools for the quality assessment of
composite indicators}.
Journal of the Royal Statistical Society, Series A, 168(2), 307--323, 2005.

\bibitem{diakoulaki1995}
D.~Diakoulaki, G.~Mavrotas, L.~Papayannakis.
\emph{Determining objective weights in multiple criteria problems: the CRITIC method}.
Computers \& Operations Research, 22(7), 763--770, 1995.

\bibitem{cherchye2007}
L.~Cherchye, W.~Moesen, N.~Rogge, T.~Van Puyenbroeck.
\emph{An introduction to `benefit of the doubt' composite indicators}.
Social Indicators Research, 82(1), 111--145, 2007.

\bibitem{mazziotta2016}
M.~Mazziotta, A.~Pareto.
\emph{On a generalized non-compensatory composite index for measuring socio-economic phenomena}.
Social Indicators Research, 127(3), 983--1003, 2016.

\bibitem{deb2002}
K.~Deb, A.~Pratap, S.~Agarwal, T.~Meyarivan.
\emph{A fast and elitist multiobjective genetic algorithm: NSGA-II}.
IEEE Transactions on Evolutionary Computation, 6(2), 182--197, 2002.

\bibitem{lahdelma2001}
R.~Lahdelma, P.~Salminen.
\emph{SMAA-2: stochastic multicriteria acceptability analysis for group decision making}.
Operations Research, 49(3), 444--454, 2001.

\bibitem{webber2010}
W.~Webber, A.~Moffat, J.~Zobel.
\emph{A similarity measure for indefinite rankings}.
ACM Transactions on Information Systems, 28(4), 1--38, 2010.

\bibitem{rousseeuw1999}
P.~J.~Rousseeuw, K.~Van Driessen.
\emph{A fast algorithm for the minimum covariance determinant estimator}.
Technometrics, 41(3), 212--223, 1999.

\bibitem{cuturi2022}
M.~Cuturi, L.~Meng-Papaxanthos, Y.~Tian, C.~Bunne, G.~Davis, O.~Teboul.
\emph{Optimal Transport Tools (OTT): a JAX toolbox for all things Wasserstein}.
arXiv:2201.12324, 2022.

\end{thebibliography}

\end{document}
